\section{Introduction}\label{sec:introduction}

Myocardial infarction (MI) damages the left ventricle unevenly. Within weeks of an occlusion, the
same ventricle contains an infarct core, a peri-infarct border region of mixed viable and
non-viable tissue, and remote myocardium that is spared by the occlusion but is not thereby normal,
since it works under altered loading. These regions do not behave alike, and the clinical questions that follow an infarction are
regional questions: which segments will recover contractile function, how far the injury extends
through the wall, and whether the remote wall is beginning to remodel. Contrast-enhanced cardiovascular magnetic
resonance (CMR) established the first of these: the transmural extent of hyperenhancement is inversely related to
the likelihood that a dysfunctional segment regains function after revascularisation~\citep{kim2000},
and infarct size measured early after primary percutaneous coronary intervention is associated with
subsequent mortality and heart failure~\citep{stone2016}. Cine-derived strain has in turn been
associated with functional recovery, adverse remodelling, remote dysfunction and subsequent
events~\citep{buss2015,reindl2021,delcanto2026,he2024}. Regional deformation is also the quantity of
interest in hypertrophic and dilated phenotypes, in cardiotoxicity surveillance and in inflammatory
and infiltrative disease, but post-infarction dysfunction is the clearest case and is used as the
running example throughout this review.

The decisions a reliable regional measurement might inform are worth naming, because they set the
accuracy that would be needed. Candidates include identifying segments likely to recover after
revascularisation, stratifying arrhythmic risk from the extent and timing of regional dysfunction,
targeting a pacing lead away from scar, and detecting early regional deterioration during
cardiotoxicity surveillance. None of these is established for cine-derived regional strain, and this
review does not claim otherwise; they are listed because a tolerance has no meaning until the
decision it serves is named.

Left ventricular ejection fraction (LVEF) is a global quantity: it summarises net chamber emptying
and therefore cannot localise a deficit that is offset by preserved or exaggerated motion
elsewhere~\citep{schulzmenger2020,hundley2022}. Global strain is a finer summary of the same kind.
In 45,700 UK Biobank participants, feature-tracking global longitudinal strain (GLS) was the only
strain measure independently associated with mortality after
adjustment~\citep{chadalavada2024}. That is useful prognostic information, but it is still a single
number for the whole ventricle. A measurement that is intended to answer a regional question has to be
evaluated as a regional measurement.

Routine balanced steady-state free-precession (bSSFP) cine CMR is the natural place to look for that
measurement. It is acquired in nearly every clinical CMR examination, needs no contrast agent, no
additional breath-holds and no dedicated pulse sequence, and large retrospective archives of it
already exist. Tagging and displacement encoding with stimulated echoes (DENSE) encode material
motion deliberately and remain the material-sensitive references, but both require a separate
acquisition, additional scan time and dedicated post-processing, and neither is routinely performed.
If regional deformation could be recovered reliably from the cine images that are already acquired,
it would be available in retrospective cohorts and in ordinary clinical practice at no extra cost in
scan time.

That possibility has driven a decade of methodological work. Commercial feature-tracking packages
brought cine strain into clinical workflows; deformable registration gave dense fields with explicit
objectives; and since 2016 a large family of learning-based estimators has appeared, first by
transplanting general-purpose unsupervised registration networks to the heart, then by designing
cardiac-specific architectures with shape, mesh, temporal, biomechanical or spectral structure, and
most recently by supervising directly on DENSE-derived displacement or on downstream tissue and
outcome labels. \Cref{sec:methods_families} reviews this literature in detail.

What has not kept pace is the evidence that these estimates recover regional abnormality. The
central question of this review is therefore:

\begin{quote}
\emph{To what extent does the available evidence support the claim that motion and strain estimated
from routine cine CMR recover regional myocardial dysfunction, meaning its magnitude, its location,
its spatial extent and its timing, and under which conditions does direct validation remain absent?}
\end{quote}

The question is deliberately symmetrical. It is not a premise that estimators smooth focal
abnormalities away, and this review does not adopt that premise; it is an empirical claim that has
been tested under some conditions and not under others, and the purpose of the funnel that follows
is to establish which is which.

Two reviews published in 2026 cover adjacent ground. Monfared et al.\ survey deep learning for
cardiac wall-motion analysis across echocardiography, four-dimensional computed tomography and cine
MRI, and emphasise data scarcity, interpretability and limited external
validation~\citep{monfared2026}. Mahmoodi et al.\ survey artificial intelligence across the whole
cardiovascular MRI pipeline, from planning and acquisition through reconstruction to analysis and
reporting~\citep{mahmoodi2026}. The present review is narrower in scope and different in
organisation: it is restricted to routine cine CMR, and it is organised by what each piece of
evidence can establish about regional recovery rather than by method family or headline performance.

\Cref{sec:observation} sets out what routine cine makes observable and what a strain number must
declare before it can be compared. \Cref{sec:methods_families} reviews the estimators.
\Cref{sec:validation} classifies the validation targets that have been used.
\Cref{sec:regional} narrows to what is known specifically about regional and focal abnormalities.
\Cref{sec:synthesis} synthesises the evidence along three axes and states what would have to be
shown for a regional attribute to be treated as recoverable. \Cref{sec:open} lists the remaining
open problems and returns to the central question.

\subsection{Review type, search and reading depth}\label{sec:methods}

This is a critical review in the sense of Grant and Booth: it evaluates and interprets a body of
evidence against an explicit question rather than aggregating effect sizes~\citep{grant2009}, and it
was written against the SANRA self-assessment items for narrative reviews~\citep{baethge2019sanra}.
Searches covered PubMed and Europe PMC, OpenAlex, Crossref and arXiv, with backward and forward
citation expansion from core method, comparison and standardisation papers, and with direct
retrieval of official dataset, repository, guideline and regulatory records. The principal search
themes were cine CMR motion estimation and deformable registration; myocardial strain measurement
and its specification; tagging, DENSE and prescribed-motion references; regional and
post-infarction strain; and accelerated or deep-learning cine reconstruction. Learning-based methods
were searched from 2016 onwards, with earlier foundational registration and tagging work included
where it is still the reference point. The evidence cutoff was 1 September 2026; one post-cutoff
study was appraised separately as a targeted update and is identified as such wherever it is
used~\citep{jahromi2026densecine}.

Two reading rules were applied throughout. First, every quantitative or comparative statement is
attributed to the specific source location it came from, and the depth at which that source was read
is recorded in the accompanying claim audit; where only an abstract was available, only
abstract-level statements are made. Second, each table cell is traceable to a source, and cells for
which the source does not report the information are marked \texttt{not reported} rather than
inferred. This review does not pool results, was screened by one person without independent
duplicate screening, was not registered and deposits no protocol, and makes no claim of exhaustive
retrieval; the supplementary search record supplies the query strings, caps and hit counts but no
screening flow counts.

\section{What cine observes and how strain is specified}\label{sec:observation}

\subsection{Information available in routine cine}\label{sec:observability}

bSSFP cine provides high blood--myocardium contrast, a well-sampled cardiac cycle and clinically
reliable ventricular geometry. At the endocardial and epicardial boundaries, moving contours and
local intensity patterns constrain motion directly. Within the myocardial wall, however, routine
cine usually lacks persistent, deliberately encoded material markers comparable with tag
intersections or DENSE phase: the trabecular and papillary structures that do provide interior
texture are neither uniformly distributed nor stable through the cycle. Correspondence between
material points is therefore inferred rather than observed, and every estimator closes the gap with
continuity, smoothness, anatomical, incompressibility, biomechanical or learned
constraints~\citep{pedrizzetti2016,duchateau2020}. This is not a property of machine learning; it
arises whenever more degrees of freedom are estimated than the image distinguishes.
\Cref{tab:signal} summarises the information each commonly paired acquisition supplies, and
\cref{fig:chain} traces the chain from image evidence to a reported regional value.

\begin{longtable}{@{}L{0.18\textwidth}L{0.23\textwidth}L{0.25\textwidth}L{0.25\textwidth}@{}}
\caption{Trackable signal and the claim supported by commonly paired CMR acquisitions. The rows are
complementary rather than an accuracy ranking.}\label{tab:signal}\\
\toprule
Acquisition & Signal that can constrain analysis & Relevant use & Principal boundary\\
\midrule
\endfirsthead
\toprule
Acquisition & Signal that can constrain analysis & Relevant use & Principal boundary\\
\midrule
\endhead
Routine bSSFP cine & Moving blood--myocardium boundaries, papillary and trabecular patterns, and incompletely persistent wall texture & Chamber volumes, boundary motion and estimator-conditioned cine strain & No deliberately encoded dense material labels; resolution, partial volume and through-plane exchange weaken regional correspondence~\citep{pedrizzetti2016,adams2023}\\
Tagging & Deforming tag lines and intersections embedded at preparation & Sparse or gridded material-sensitive displacement and strain comparison & Tag fading, spatial resolution, tag-tracking and registration error; a dedicated acquisition~\citep{cao2018,wehner2018}\\
DENSE & Motion-encoded phase that yields displacement components under its acquisition model & Material-sensitive displacement and strain reference with explicit uncertainty & Phase unwrapping, stimulated-echo acquisition, spatial and temporal registration and processing error; not an error-free truth~\citep{cao2018,wang2023strainnet}\\
Late gadolinium enhancement & Contrast-related tissue signal delineating infarct core, peri-infarct and remote zones & Independent tissue zoning and biological association & Different breath-holds, cardiac phases and slice geometry require registration; tissue contrast does not encode dense motion~\citep{guo2021,almeida2021}\\
\bottomrule
\end{longtable}

\begin{figure}[!ht]
 \centering
 \includegraphics[width=\textwidth,keepaspectratio]{figures/p1_Figure1_image_information_to_strain_report.pdf}
 \caption{From image information to a regional strain report. \textbf{a}, end-diastolic (ED) and
 end-systolic (ES) cine frames supply boundaries, texture and intensity, but not labels identifying
 the same material point at both phases. \textbf{b}, an estimator infers a candidate correspondence
 $\vect{\varphi}(\vect{X},t)$; feature-tracking, registration, biomechanical and learning-based
 families overlap and each constrains the undetermined part of the field through its own
 assumptions. \textbf{c}, strain is then defined on that inferred mapping through the displacement
 $\vect{u}=\vect{\varphi}-\vect{X}$, the deformation gradient $\mat{F}$ and, in this example, the
 Green--Lagrange tensor $\mat{E}$; seven further specification choices are required. \textbf{d}, the
 regional value that is finally reported is conditional on all of the preceding choices. The cine
 frames in panels a and c are real images; the ring, the dashed correspondences and the strain curve
 are schematics and carry no measured value, which is why the curve is labelled illustrative and
 normalised rather than given units.}
 \label{fig:chain}
\end{figure}

\subsection{Through-plane motion, multiple views and image quality}\label{sec:throughplane}

Two-dimensional regional analysis is particularly exposed to through-plane motion. A short-axis
pixel at one phase need not contain the same tissue at the next: basal longitudinal excursion, slice
thickness and non-parallel repeat prescriptions all change the sampled material, and a
two-dimensional estimator can only explain that change as in-plane deformation. Global averaging
absorbs non-systematic local error, whereas peak and time-to-peak estimates remain sensitive to it.
Consistent with this observation model, the same-day strain-8 experiment found substantially poorer
repeatability for many segmental longitudinal and radial measures than for global measures, with
basal segments the most difficult~\citep{bell2025strain8}; the pattern is consistent with
through-plane exchange but does not attribute every error to it.

Long-axis cine narrows part of the gap by constraining longitudinal excursion and apical--basal
geometry that a short-axis stack represents poorly, and multi-view and three-dimensional estimators
exploit exactly this complementary information~\citep{meng2022mulvimotion,meng2024deepmesh,liu2026fe}.
Independently acquired two-dimensional slices are not a simultaneous isotropic material volume,
however, and slice misregistration remains possible. Additional views constrain the field without
guaranteeing identifiable material correspondence, so a method claiming three-dimensional motion
from multi-planar two-dimensional cine should report slice geometry, temporal alignment, view
interaction and behaviour when one view is missing or degraded~\citep{duchateau2020}.

Image quality enters the same way. Temporal undersampling can alter peak magnitude and timing;
spatial averaging can alter the width and amplitude of a regional gradient; mistriggering and
arrhythmia disrupt the assumed cycle; and contour correction changes the information supplied to the
tracker. Silently excluding low-quality cases may overstate accuracy for a clinical population,
while retaining them without a failure flag hides the resulting uncertainty, so a validation dataset
needs both acceptable and deliberately challenging cases with the rejection rate and the residual
error reported together. \Cref{sec:conditions} returns to this with the reconstruction literature.

\subsection{Tissue zones after infarction}\label{sec:zones}

Because post-infarction dysfunction is this review's running example, the definition of these
regions needs to be precise. Late gadolinium enhancement (LGE) separates infarct core,
peri-infarct border and remote myocardium by contrast kinetics; the border or grey zone and the
extent of peri-infarct scar have each been related to arrhythmic and clinical outcomes in their own
cohorts~\citep{yang2013,tulumen2021}. These zones are defined by \emph{tissue} signal. Many of
the comparisons discussed below are scar-versus-remote within the same heart, which is an efficient
design because it removes between-subject variation; remote myocardium after infarction is, however, a
loaded, remodelling region rather than a healthy control, so a scar-to-remote contrast understates
the difference from a normal ventricle by an unknown amount. These zones are an
independent and clinically meaningful stratification, and they are the right way to ask whether a
strain measurement differs between injured and uninjured myocardium. They are not, however, a motion
reference, and an LGE map does not state how far any material point moved. The distinction governs
how the evidence in \cref{sec:cineonly} may be read.

Using the two together adds an alignment problem of its own. In the porcine registration study of
Guo et al., cine and LGE were acquired with different contrasts and required explicit registration
before infarct heterogeneity could be quantified~\citep{guo2021}. Even successful anatomical
registration does not establish that the cine tracker recovered material motion within each zone, so
registration uncertainty should be propagated into zone-wise strain analyses, or at least probed
through plausible perturbations of slice, phase and contour alignment.

\subsection{Specifying the strain measurement}\label{sec:estimand}

A strain value is not defined by the word \enquote{strain}. Consider a forward material mapping
$\vect{x}=\vect{\varphi}(\vect{X},t)$ from a reference coordinate $\vect{X}$ to the current
coordinate $\vect{x}$, with displacement $\vect{u}(\vect{X},t)=\vect{\varphi}(\vect{X},t)-\vect{X}$
and deformation gradient
$\mat{F}=\partial\vect{\varphi}/\partial\vect{X}=\mat{I}+\partial\vect{u}/\partial\vect{X}$, where
$\mat{I}$ is the identity tensor. Writing $\mat{C}=\mat{F}^{\trans}\mat{F}$ and
$\mat{B}=\mat{F}\mat{F}^{\trans}$ for the right and left Cauchy--Green tensors, the Green--Lagrange
strain $\mat{E}=\tfrac12(\mat{C}-\mat{I})$ is referred to the reference configuration and the
Euler--Almansi strain $\mat{e}=\tfrac12(\mat{I}-\mat{B}^{-1})$ to the current one. For a stretch
$\lambda_i$ in principal direction $i$, engineering strain is $\lambda_i-1$, Green--Lagrange strain
is $(\lambda_i^2-1)/2$, Euler--Almansi strain is $(1-\lambda_i^{-2})/2$, and Hencky or logarithmic
strain is $\ln\lambda_i$. At $\lambda=0.8$ these give $-0.20$, $-0.18$, $\approx-0.281$ and
$\approx-0.223$ respectively: the same shortening, four non-interchangeable percentages. The
calculation follows from the definitions and is not a claim about any physiological value, but it
sets the scale of the problem, because differences of this size are comparable with the between-method
differences reported in \cref{sec:crosssoftware}.

Standard terminology from deformation imaging provides a common language for these
quantities~\citep{voigt2015,amzulescu2019,smiseth2025}, but a shared vocabulary does not prove that
a given CMR implementation follows every convention, and in clinical usage the word
\enquote{Lagrangian} can denote a relative length change rather than a Green--Lagrange
tensor~\citep{voigt2015}. Several further choices change the number as much as the tensor does. Differentiating a backward
sampling map as though it were the forward material map changes the tensor and can change signs and
directions, and negating a displacement field does not in general invert a nonlinear transformation.
The reference phase, whether end-diastole, the first acquired frame, end-systole or a user selection,
sets the baseline length. Circumferential, radial and longitudinal axes may be built from a cavity
centre, a local wall normal, a fitted surface or a ventricular mesh. Layer-specific strain may mean
a contour length, a band average, a sampled subendocardial surface or a transmural interpolation,
and radial strain is the most heterogeneous of all: a tensor projection, wall thickening, a change
in endocardial--epicardial distance or an average of local radial components. Finally, the peak of
an averaged curve is generally not the average of local peaks, because regions reach their extrema
at different phases, so taking a peak and averaging across regions do not commute.
\Cref{tab:estimand} is the minimum that must be declared before two values may be compared.

The comparison by Jahromi et al.\ shows how much this matters even between careful studies:
single-slice cine strain was built from contour-length, wall-thickness and myocardial-area ratios
with longitudinal shortening inferred under uniform-shortening and per-slice volume-conservation
assumptions, whereas two-slice DENSE used displacement gradients. The two shared a strain measure
and a component label without sharing material sampling or spatial
aggregation~\citep[Eqs.~1--6]{jahromi2026densecine}. Conventions can be declared explicitly: the
United States Food and Drug Administration record K232661 states that one cleared product uses an
end-diastole-referenced Lagrangian tensor in radial, circumferential and longitudinal directions and
defines peak as the maximum absolute value over the cycle~\citep{fdaK232661}, which is exactly the
kind of declaration \cref{tab:estimand} asks for, though it licenses no assumption about other
products or other releases.

\begin{longtable}{@{}L{0.17\textwidth}L{0.26\textwidth}L{0.24\textwidth}L{0.24\textwidth}@{}}
\caption{Minimum strain measurement specification. Two values are not safely compared unless every
row is compatible or a conversion is justified.}\label{tab:estimand}\\
\toprule
Specification item & Required declaration & Why it changes the value & Minimum audit question\\
\midrule
\endfirsthead
\toprule
Specification item & Required declaration & Why it changes the value & Minimum audit question\\
\midrule
\endhead
Mapping and reference & Forward material map or backward spatial warp, composition order, and the reference phase (end-diastole, first frame, end-systole or other) & Inverse and gradient operations do not commute, and the reference sets the baseline length & Which coordinates does each vector index, was the map inverted before differentiation, and is the reference identical across subjects?\\
Strain measure & Engineering, Green--Lagrange, Euler--Almansi, Hencky/logarithmic or other & At finite deformation the measures differ numerically from the same mapping & Is the tensor or formula stated, rather than inferred from a label?\\
Direction & Local radial, circumferential and longitudinal axes, their origin and their update rule & Rotating bases and different centre definitions change the projection & Are the axes material, spatial, slice-based or three-dimensional anatomical?\\
Spatial object & Endocardial contour, epicardial contour, midwall, named layer, voxel band or full myocardium & A contour length is not a dense transmural tensor average & What tissue or geometric object contributes to the reported scalar?\\
Spatial differentiation & Voxel-, mesh- or basis-based deformation gradient; interpolation order; derivative operator; spatial smoothing; boundary handling & Strain is a spatial derivative, so differentiation and smoothing can amplify or attenuate regional displacement error & How was the deformation gradient computed, regularised and evaluated near the myocardial boundaries?\\
Temporal summary and aggregation & End-systolic value, peak of the global curve, mean of segmental peaks or peak strain rate; area, length or volume weighting; American Heart Association (AHA) segmentation and excluded segments; units and sign convention & Peak and averaging do not commute, frame rate affects extrema, and weighting and missing segments alter both global and regional summaries & Was the extremum observed at an acquired frame or interpolated, and can signed curves, timestamps, weights and exclusion flags be exported?\\
\bottomrule
\end{longtable}

Finally, strain describes deformation. It is not stress, active tension or contractility; inferring
those requires assumptions about loading, geometry, constitutive behaviour and activation, and a
homogeneous material model may not represent both viable and scarred tissue~\citep{almeida2021,amzulescu2019}.
Radial thickening in particular can reflect kinematic coupling with longitudinal and circumferential
shortening under near-volume conservation rather than a locally generated radial force. This review
treats strain as an image-derived kinematic measurement throughout.

\section{Motion estimation from cine: from conventional to learning-based methods}\label{sec:methods_families}

\subsection{An underdetermined problem, and what the prior is doing}\label{sec:illposed}

Let $I_t(\vect{x})$ denote image intensity at spatial coordinate $\vect{x}$ in phase $t$, and let
$\vect{\varphi}_{0\rightarrow t}$ map a material point in a reference configuration to its position
at $t$. Almost every estimator reviewed below can be written as, or interpreted through, the
objective
\begin{equation}
\widehat{\vect{\varphi}}_{0\rightarrow t}
=\arg\min_{\vect{\varphi}_{0\rightarrow t}}
\mathcal{D}\!\left(I_0,I_t\circ\vect{\varphi}_{0\rightarrow t}\right)
+\alpha\,\mathcal{R}(\vect{\varphi}_{0\rightarrow t}),
\label{eq:registration}
\end{equation}
in which $\mathcal{D}$ rewards image agreement, $\mathcal{R}$ encodes prior structure and $\alpha$
sets their balance. Here $\vect{\varphi}_{0\rightarrow t}$ is a forward map defined on the reference
domain, so that the target image is evaluated as $I_t(\vect{\varphi}_{0\rightarrow t}(\vect{X}))$;
frameworks that optimise a backward sampling map use a different coordinate convention, and the two
must not be interchanged when transformations are composed or differentiated for strain.
\Cref{eq:registration} is written for one pair of phases; the groupwise and sequence models in
\cref{sec:conventional,sec:cardiacspecific} optimise an analogous objective over the whole cycle at
once, and learning-based methods minimise its expectation over a training set rather than solving it
per case.

The difficulty is that $\mathcal{D}$ alone does not identify the map. Where intramyocardial texture
is weak, many fields achieve nearly the same data term while implying different displacement
gradients, and the prior selects among them. \Cref{fig:correspondence} makes the point exactly: for
a circular annulus, the map $\varphi(r,\theta)=(r,\theta+a\sin\theta)$ with $|a|<1$ has positive
angular derivative $1+a\cos\theta$, so it is a smooth one-to-one reparameterisation that leaves both
boundaries and the continuous binary mask unchanged. Dice is 1 and the Hausdorff distance is 0,
while the tangential stretch varies between $1-|a|$ and $1+|a|$. The construction is analytic, not
an empirical error estimate, but it establishes that image and contour agreement cannot by
themselves certify material correspondence, which is why \cref{sec:validation} treats overlap and
material-sensitive references as different kinds of evidence.

\begin{figure}[!t]
 \centering
 \includegraphics[width=\textwidth,keepaspectratio]{figures/p2_Figure2_same_contours_not_same_correspondence.pdf}
 \caption{Identical contours do not determine the same material correspondence. \textbf{a}, routine
 cine supplies anatomy and image cues, tagging supplies a prepared magnetisation pattern, and
 displacement encoding with stimulated echoes (DENSE) encodes displacement in the magnetic resonance
 phase. \textbf{b}, under the analytic map $\theta'=\theta+a\sin\theta$ with $a=0.5$, eight labelled
 material points are redistributed around the annulus while the annulus itself is unchanged.
 \textbf{c}, the identical continuous masks give Dice~$=1$ and Hausdorff distance~$=0$, whereas the
 tangential stretch $\lambda_\theta=1+a\cos\theta$ ranges from 0.5 to 1.5. The stretch ratio is not
 Green--Lagrange strain, and discrete resampling need not preserve exact equality. This is a derived
 counterexample, not a measurement of any estimator's error. Panel a reproduces real acquired
 images, credited within the figure.}
 \label{fig:correspondence}
\end{figure}

Two statements about priors are easy to conflate and are kept separate throughout this review.
The first concerns what a representation can express. Strain is a spatial \emph{derivative} of
displacement, so a displacement field that looks smooth to the eye can still carry a rapidly varying
strain; visual smoothness of $\vect{u}$ is not evidence that $\partial\vect{u}/\partial\vect{X}$ has
been flattened. Restricting a displacement field to a band of spatial frequencies does restrict the
derivatives it can represent, but the consequence for a particular regional measurement depends
jointly on the bandwidth, the spatial scale of the abnormality and the measurement definition, and a
Fourier basis is global: widening the band raises the representable frequency content everywhere and
does not target a chosen region. The second statement concerns the estimation process: even with a
representation rich enough to express a focal deficit, the optimiser or network may not recover it
if the image evidence is weak or the penalty $\alpha\mathcal{R}$ is strong. A third factor, the
information actually present in the input images, is separate again.

The three come apart in practice, and the way they come apart is the basis of an attribution
experiment. Suppose a focal deficit is under-recovered. If the same estimator recovers it when the
penalty weight $\alpha$ is reduced, the limitation was the estimation process. If it is still
under-recovered after the deficit has been projected directly onto the estimator's function space
(a cheap test that uses the known field rather than the images), the limitation lies in the
representation. And if neither holds, but the deficit becomes recoverable when the same motion is
re-imaged at finer resolution or with added long-axis views, the limitation was the input
information. Attributing an observed regional error therefore requires an experiment that holds the
other two factors fixed; this is the structure of \cref{sec:synthesis}.

Temporal strategy adds a further set of assumptions. Reference-to-frame tracking keeps one reference
configuration but must estimate large displacements from it; sequential tracking estimates smaller
inter-frame steps but accumulates composition error and drift; groupwise methods estimate all phases
jointly and can improve temporal consistency under whatever temporal constraint they
impose~\citep{metz2011,chen2024}. A periodic deformation can close perfectly at the end of the cycle
and still be wrong at intermediate phases, so cycle closure is a useful constraint and a useful
quality-control signal rather than proof of correct material trajectories.

\subsection{Conventional estimators}\label{sec:conventional}

A note on evaluation applies to this whole section. Segmentation does not merely precede motion
estimation; it can initialise the search region, define the cardiac coordinate system, select the
sampled layer and, in several learning-based designs, supply an anatomical loss during training.
When the same reference contours both shape the fitted transformation and provide the overlap score
that assesses it, the evaluation is not independent of the fitting. A single inter-observer
coefficient for the final scalar cannot separate contour uncertainty from tracking uncertainty from
strain-computation uncertainty; perturbing each in turn can.

Boundary feature tracking follows local patterns near user-defined or automatically generated
contours, usually with propagation and correction. It exploits the strongest signal routine cine
provides and keeps contour interaction explicit, but interior motion is to a large degree
interpolated from boundary behaviour~\citep{pedrizzetti2016}. It is the family implemented in the
commercial packages whose mutual agreement is examined in \cref{sec:crosssoftware}. Automating its
front end was an early application of deep learning: Vigneault et al.\ segmented the sequence with a
fully convolutional network and then fitted quadratic basis splines simultaneously across all cardiac
frames, reporting pixel accuracy around 97\% against expert manual segmentation in 42 patients with
hypertrophic cardiomyopathy and 21 controls~\citep{vigneault2017}, and Puyol-Ant\'on et al.\
combined convolutional-network segmentation with subsequent feature tracking and strain
calculation~\citep{puyolanton2018}. In both, the learned component is the segmentation and the
motion model remains a conventional tracker.

Intensity-based deformable registration estimates a dense field by optimising
\cref{eq:registration} directly. The free-form deformation model of Rueckert et al., a B-spline
control lattice with a smoothness penalty, remains the reference implementation of this
approach~\citep{rueckert1999}. Groupwise variants remove the dependence on an arbitrarily chosen
reference frame: Metz et al.\ fit an $n$D$+t$ B-spline model to the whole dynamic sequence with
spatial and temporal smoothness and an optional cyclic constraint, reporting sub-voxel accuracy on
synthetic and thoracic computed tomography data~\citep{metz2011}. More recently, Chen et al.\
replaced the groupwise dissimilarity term with a locally low-rank metric, evaluating it on a
20-case simulation, 100 public cine cases and 16 in-house tagging-MRI patients, and reporting
reduced late-diastolic drift and better correlation with tagging for radial strain in several
myocardial segments and levels than optical flow, pairwise registration, global low-rank groupwise
registration, $n$D$+t$ B-splines, and pTVreg, a general-purpose total-variation-regularised
groupwise registration package~\citep{chen2024}.

Mechanical constraints narrow the admissible set in a different way. Nearly incompressible models
impose a volume constraint whose adequacy depends on the conditions studied; the FDA record K232661
describes one cleared implementation as a two-dimensional nearly incompressible deformable model
with control points on a midwall curve, without disclosing every numerical
detail~\citep{fdaK232661}. Geometry- and anatomy-constrained estimators occupy similar ground,
constraining the field with masks, contours, shape priors or geometric correspondence rather than a
constitutive model~\citep{chen2020anatomy,hammouda2020}.

Tagging analysis is the comparison point rather than a cine method. Harmonic phase (HARP) analysis
extracts displacement from the phase of an isolated spectral peak of an image tagged by spatial
modulation of magnetisation (SPAMM), giving material-sensitive
motion directly from the acquisition~\citep{osman2000harp}. The contrast is instructive: HARP solves
a different problem, because the information it needs has been deliberately encoded into the images.

\subsection{Unsupervised learning-based registration}\label{sec:unsupervised}

The learning-based literature began by amortising \cref{eq:registration}. VoxelMorph reformulated
registration as a function, parameterised by a convolutional network, that maps an image pair to a
deformation field; the unsupervised training objective is the same image-matching and smoothness
objective that a classical method optimises per pair, and the gain is that inference becomes a
single forward pass~\citep{balakrishnan2019voxelmorph}. TransMorph replaced the convolutional
encoder with a hybrid Transformer--convolutional architecture to widen the receptive field, and
added diffeomorphic and Bayesian variants that respectively guarantee topology preservation and
produce an uncertainty estimate; its reported validation covers inter-patient and atlas-to-patient
brain MRI and phantom-to-computed-tomography registration~\citep{chen2022transmorph}. The important
point for this review is not the backbone but the objective: in this family the prior is still
smoothness or diffeomorphism, now partly learned, and the reported evidence is label overlap and
deformation-field regularity.

A distinct strand constrains the representation rather than the penalty. Fourier-Net predicts a
band-limited representation of the displacement field in a low-dimensional frequency domain and
recovers the full-resolution field by zero-padding and an inverse discrete Fourier transform,
replacing the expensive full-resolution decoder of a U-Net-style registration network and
substantially reducing computation~\citep{jia2023fouriernet}; Fourier-Net+ extends the idea by
taking band-limited representations of the images themselves as input~\citep{jia2026fouriernetplus}.
Band-limiting is an explicit, tunable statement about which deformations the model may output, which
makes this family unusually well suited to the experiment described in \cref{sec:axis_prior};
but, as \cref{sec:illposed} sets out, the bandwidth is a global property of the field, and what it
implies for a focal regional strain has to be measured rather than asserted. Both are general-purpose registration
networks rather than cardiac methods: the evaluations reported in the records retrieved here use
general medical image registration benchmarks with overlap and deformation-regularity metrics, not
cine CMR against a motion reference. Within the searches recorded in \cref{sec:methods} and up to
the stated cutoff, this review did not identify an evaluation of a band-limited cine estimator against a
regional motion or strain reference; that is a statement about what those searches returned, not an
assertion that no such study exists.

\subsection{Cardiac-specific motion estimation}\label{sec:cardiacspecific}

A second group of methods builds cardiac structure into the estimator. Qin et al.\ learned motion
estimation and segmentation jointly in a single architecture, an unsupervised Siamese recurrent
spatial-transformer branch for motion and a fully convolutional branch for segmentation, sharing a
multi-scale feature encoder so that each task supplies features to the other, and reported
improved accuracy and speed over separate models in 220 subjects~\citep{qin2018joint}. Zheng et al.\ learned a pixel-wise
\enquote{apparent flow} between two phases of a 2D$+t$ cine sequence semi-supervised from a large
unsegmented set and a small expert-segmented set, then combined the flow with segmentation masks to
derive radius and wall-thickness time series per segment; using nine feature values, their
classifier reached 95\% and 94\% accuracy on the ACDC training and testing
sets~\citep{zheng2019apparent}. The target in that study is pathology classification, so the motion
representation is validated through its usefulness for a downstream label rather than against a
motion reference, a pattern that recurs in \cref{sec:cineonly}. Motion Pyramid Networks introduced
multi-scale fused estimation with a cyclic teacher--student training procedure, which is a training
strategy and not evidence of cardiac-cycle closure~\citep{yu2020mpn}, and FOAL added online
adaptation of a learned estimator to a new sequence at test time~\citep{yu2020foal}.

Multi-view and mesh-based designs target the through-plane problem of \cref{sec:throughplane}.
MulViMotion fuses short-axis and long-axis cine in a hybrid 2D/3D network and adds a shape
regularisation module that provides weak supervision to the 3D field, evaluated on 580 UK Biobank
subjects~\citep{meng2022mulvimotion}. DeepMesh represents the heart as a 3D mesh of endocardial and
epicardial surfaces, propagates a template mesh into subject space and estimates vertex-wise
displacement through a differentiable mesh-to-image rasteriser, so that vertex correspondence is
maintained across frames and subjects~\citep{meng2024deepmesh}. CardioMorphNet takes the shape idea
further: a recurrent variational autoencoder with separate posteriors for bi-ventricular
segmentation and for motion, trained by recursively registering segmentation maps \emph{without} an
intensity-based similarity loss, and validated on UK Biobank and M\&Ms by comparing warped masks
with ground-truth masks; its Bayesian formulation also yields uncertainty maps for the estimated
field~\citep{movahed2026cardiomorphnet}. Qiao et al.\ take a different route to the same problem,
representing cardiac motion through keypoints extracted without supervision and reconstructing the
dense field as a weighted combination of basis vectors derived from matched keypoint pairs, with
evaluation on four-dimensional computed tomography and cine MRI and a strain analysis reported to
separate patients under different cardiac conditions~\citep{qiao2026stcmt}.

Temporal structure has been addressed in parallel. A
low-rank groupwise deformation model has been proposed for cine tracking in a preprint that is a
master's thesis and has not been peer reviewed~\citep{rendell2024lowrank}, and a Mamba-based sequential
model has been presented in a MICCAI workshop paper~\citep{yin2025mcm}. Ruan et al.\ embed the
deformation in a continuous-time latent neural ordinary differential equation with bidirectional
forward--backward evolution and a bidirectional Lagrangian consistency term, evaluated on ACDC and
M\&Ms~\citep{ruan2026latentode}.

One parallel development belongs here precisely because it is not a motion-estimation development.
CMR-specific foundation models now exist: one was pre-trained on 15 million cine images from 74,916
studies and fine-tuned for segmentation, landmark localisation, disease diagnosis and
prognostication across eight independent datasets~\citep{fu2026cinema}, and another was trained by
self-distillation on 36 million images from 27,524 subjects and fine-tuned for nine tasks including
classification, segmentation, landmark localisation and pathology detection~\citep{jacob2025cmrfm}.
Dense motion estimation is not among the tasks either reports. That is informative about where
effort is currently going, and it means the pre-trained representations these models offer have not
yet been tested on the problem this review is about.

Read together, this group has made real progress on two things the earlier literature handled badly:
three-dimensional motion from two-dimensional acquisitions, and temporal consistency across the
cycle. Their reported validation needs stating plainly, because it bears
directly on this review's question. In the cardiac-specific papers reviewed here, the dominant
evidence is agreement between a warped segmentation and a reference segmentation, together with
surface distances and deformation-field regularity; landmark error appears where a landmarked
dataset was available. That is appropriate evidence for the claims those papers make, and it is also
the evidence that \cref{fig:correspondence} shows cannot by itself establish material
correspondence.

\subsection{Physics, biomechanics and continuous representations}\label{sec:physics}

A third group replaces a generic smoothness penalty with a statement about tissue mechanics or about
the function space itself. Qin et al.\ learned the prior instead of writing it down: a variational
autoencoder trained on biomechanically simulated deformations defines a manifold of plausible
fields, and tracking becomes a search over that manifold, validated on three public cine
datasets~\citep{qin2023biomech}. WarpPINN takes the opposite route and states the physics
explicitly, representing the deformation with a neural network trained per case under an
intensity-similarity term plus a hyperelastic regulariser and a Jacobian penalty for
near-incompressibility; strain follows by automatic differentiation, and Fourier feature mappings
are used to overcome the spectral bias of the network and so to allow sharp features in the strain
field. It was tested on synthetic examples and on a cine bSSFP benchmark of 15 healthy
volunteers~\citep{arratia2023warppinn}; a later version adds a fibre-stretch term and reports
evaluation on a synthetic phantom and the same 15-participant healthy
benchmark~\citep{alvarez2026warppinn}. Implicit neural representations of intramyocardial motion and strain
have also been presented in the 2026 STACOM proceedings; no abstract accompanies the proceedings
record retrieved for this review, so the acquisition, data and evaluation of that study are not
characterised here~\citep{bell2026inr}.

Liu et al.\ combine the mechanical and the temporal arguments in a space-time finite-element
digital image correlation framework with region-specific biomechanical regularisation, a data-driven
temporal decomposition and a correlation-based multi-view alignment module. Their evaluation is the
most layered in this group: one synthetic dataset with ground-truth motion and strain fields, three
public datasets with ground-truth landmarks or myocardial masks, and a clinical dataset with
ground-truth myocardial masks, compared against two classical feature-tracking methods and four
deep-learning baselines~\citep{liu2026fe}. The limiting factor is the same throughout the group:
where a strain reference exists at all, it exists in one simulation setting, and the in-vivo
evaluation returns to landmarks and masks. A mechanical prior is not a weaker constraint than a
smoothness prior, merely a differently shaped one, and whether it preserves a focal deficit depends
on whether its constitutive assumptions hold in heterogeneous injured tissue, which is an open
question rather than a defect.

\subsection{Reference-supervised and end-to-end strain}\label{sec:supervised}

The final group changes the supervision rather than the prior. DeepTag estimates motion on tagged
images without labels, using a bidirectional diffeomorphic registration network between consecutive
frames and a differentiable composition layer to recover the Lagrangian field, validated on a
clinical tagging dataset by landmark tracking accuracy~\citep{ye2021deeptag}. StrainNet is the
clearest example of transferring material-sensitive information into cine analysis. A 2D$+t$
convolutional network was trained to predict intramyocardial displacement from the motion of
myocardial contours, with DENSE displacement measurements as ground truth, across 161 patients, 99
healthy adults and 45 children and adolescents. It achieved an end-point error of
$0.75\pm0.35$~mm against DENSE, and when applied to contour motion from cine it reached intraclass
correlation coefficients (ICCs) with DENSE of 0.87 for global and 0.75 for segmental circumferential
strain, against 0.72 and 0.48 for commercial feature tracking~\citep{wang2023strainnet}.

The same supervision route has since been extended to the part of the motion that a displacement
network handles least naturally. Lei et al.\ report a DENSE-supervised network for routine cine
bSSFP that estimates the large rotational dynamics of the myocardium in a separate time-series
branch, predicts radial motion by deformable registration in a second branch and fuses the two,
motivated by the observation that earlier DENSE-supervised models captured twist and torsion
poorly~\citep{lei2025denserotation}. A pre-clinical study takes the idea to a different reference:
MyoNet was trained to recover displacement maps from cine images in a rat model of
radiation-therapy-induced injury and evaluated against strain derived by a sine-wave modelling
method, reporting intraclass correlations of 0.973 and 0.975 for circumferential and radial strain
respectively~\citep{an2026myonet}. Both are abstract-level readings here. Neither changes the
validation pattern set out in \cref{sec:method_table}, and the second is reported in a rat model
against an analysis-derived reference, so it bears on the design rather than on human cine accuracy.

DeepStrain is the most completely documented end-to-end cine workflow. It centres the ventricle,
segments anatomy with CarSON, estimates end-diastole-to-phase displacement with CarMEN and derives
Green--Lagrange radial and circumferential strain. The final motion model was trained on all 150
ACDC subjects, including the 50-subject official segmentation test set with predicted labels, and
the original study reported a Cardiac Motion Analysis Challenge (CMAC) landmark end-point error of
$2.9\pm1.5$~mm together with scan--rescan repeatability in ten healthy
volunteers~\citep{morales2021}. A later single-centre comparison in 47 healthy participants and 533
patients found correlations with manual cvi42 analysis of 0.83--0.85 for global radial and
0.88--0.91 for global circumferential strain, with mean processing times of $4.1\pm1.2$~minutes for
manual feature tracking and $34.7\pm3.3$~seconds for DeepStrain~\citep{morales2022comparison}; the
comparator is another estimator rather than a motion reference, and the developer overlap limits
independence. The pinned public repository supplies inference code, examples and an MIT licence, but not a
complete deterministic retraining workflow or a supported clinical release~\citep{deepstrainrepo}.
An unrelated tagging network shares the name~\citep{abdelmoniem2021}.

A further set of systems is supervised on a downstream label instead of on motion. Deep learning
synthetic strain derives regional wall-motion assessment from cine and was evaluated against the
consensus of four cardiothoracic radiologists in 53 patients and 846 segments, reaching an area under
the receiver operating characteristic curve (AUC) of 0.90 and sensitivity, specificity and accuracy
of 86\%, 85\% and 86\% at a 30\% peak radial strain threshold~\citep{masutani2023dlss}. Deep
learning-based aligned strain was developed to detect fibrotic myocardium in Duchenne muscular
dystrophy, with LGE as the reference~\citep{koehler2025alignedstrain}. Bello et al.\ supervised a
motion representation directly on survival, reporting a Harrell C index of 0.75 (95\% confidence
interval 0.70--0.79) against a human benchmark of 0.59 (0.53--0.65) in 302
patients~\citep{bello2019}. Each is a well-posed study, and each validates something other than the
accuracy of the motion field: the supervision target determines what the evidence can support, and
\cref{sec:cineonly} returns to this distinction where it matters most.

\subsection{Comparison across methods}\label{sec:method_table}

\Cref{tab:methods_what} records what each method does and \cref{tab:methods_validation} records how
it has been validated, in the same row order. Every cell is drawn from the cited source; where the
source does not report an item, the cell reads \texttt{not reported}, and no value has been inferred
from a related paper. \Cref{tab:methods_validation} is the one that bears on this review's question,
and reading its last column downwards is the quickest way to see where the evidence thins. The rows
are not ranked and are not equivalent in status: they include peer-reviewed journal and conference
papers, a workshop paper, a preprint submitted as a master's thesis and one regulatory summary, and
the \enquote{Year, source} column states which is which. A method with a short validation entry has
usually been evaluated for the claim its own authors made, not tested and found wanting for the
claim this review is asking about.

\begin{longtable}{@{}L{0.135\textwidth}L{0.115\textwidth}L{0.135\textwidth}L{0.155\textwidth}L{0.215\textwidth}L{0.155\textwidth}@{}}
\caption{What each estimator does. Entries are taken from the cited source; \texttt{not reported}
means the source does not state the item. Hybrid systems can belong to more than one group.}\label{tab:methods_what}\\
\toprule
Method & Year, source & Input & Supervision & Prior or representation & Output\\
\midrule
\endfirsthead
\toprule
Method & Year, source & Input & Supervision & Prior or representation & Output\\
\midrule
\endhead
\multicolumn{6}{@{}l}{\textit{Conventional estimators}}\\
Boundary feature tracking & Family; see \citep{pedrizzetti2016} & Cine plus endocardial and epicardial contours & None (per-case tracking) & Contour continuity and propagation; interior motion largely interpolated from boundaries & Boundary and segmental displacement, velocity and strain\\
Automated segmentation plus spline tracking & 2017, FIMH~\citep{vigneault2017}; 2018, ISBI~\citep{puyolanton2018} & Cine sequence & Supervised segmentation; tracking itself unsupervised & Learned segmentation front end feeding a conventional spline or feature tracker & Contours, then boundary-derived strain\\
Free-form deformation (FFD) & 1999, IEEE TMI~\citep{rueckert1999} & Image pair & None (per-pair optimisation) & Affine plus B-spline control lattice with a smoothness cost; normalised mutual information data term & Dense displacement field\\
Groupwise $n$D$+t$ B-splines & 2011, Med Image Anal~\citep{metz2011} & Whole dynamic sequence & None & $n$D$+t$ B-spline with spatial and temporal smoothness; optional cyclic constraint; no fixed reference frame & Dense 2D/3D$+t$ displacement\\
Groupwise locally low-rank (Groupwise-LLR) & 2024, BMC Med Imaging~\citep{chen2024} & Whole cine sequence & None & Groupwise registration minimising the sum of patchwise signal ranks of the deformed movie & Whole-cycle displacement field and strain\\
Anatomy-Aware Tracker & 2020, MLMI~\citep{chen2020anatomy} & Long-axis cine & Weak, from a shape model & Convolutional variational autoencoder encoding realistic myocardial shapes, used to refine a dense tracker & Anatomy-aware dense motion field\\
Laplace-based murine framework & 2020, Sci Rep~\citep{hammouda2020} & Mouse cine & Supervised segmentation networks & Laplace equation solved between contours of successive frames & LV wall point tracks and Lagrangian strain\\
Cleared nearly in\-com\-pressible model & 2023, FDA K232661 \citep{fdaK232661} & bSSFP cine with confirmed contours & not reported & Two-dimensional nearly incompressible deformable model with control points on a midwall curve & Global and regional strain; end-diastole-referenced Lagrangian tensor\\
Harmonic phase (HARP) & 2000, IEEE TMI~\citep{osman2000harp} & SPAMM-tagged MR (comparator, not cine) & None & Isolated spectral peak of the tagged image; phase of the complex image & Displacement for small motions and 2D strain\\
\midrule
\multicolumn{6}{@{}l}{\textit{Unsupervised learning-based registration}}\\
VoxelMorph & 2019, IEEE TMI~\citep{balakrishnan2019voxelmorph} & Image pair & Unsupervised image matching; optionally auxiliary segmentations & Learned amortisation of a similarity-plus-smoothness objective & Dense deformation field in one forward pass\\
TransMorph & 2022, Med Image Anal~\citep{chen2022transmorph} & Volumetric image pair & Unsupervised & Hybrid Transformer--convolutional encoder; diffeomorphic and Bayesian variants & Deformation field; uncertainty in the Bayesian variant\\
Fourier-Net & 2023, AAAI~\citep{jia2023fouriernet} & Image pair & Unsupervised & Low-dimensional band-limited Fourier representation of the displacement, expanded by a parameter-free model-driven decoder & Full-resolution displacement field\\
Fourier-Net+ & 2026, IEEE TNNLS~\citep{jia2026fouriernetplus} & Band-limited spatial representation of the image pair & Unsupervised & As Fourier-Net, with a shortened contracting path; a cascaded variant is also proposed & Full-resolution displacement field\\
\midrule
\multicolumn{6}{@{}l}{\textit{Cardiac-specific estimators}}\\
Joint motion and segmentation & 2018, MICCAI~\citep{qin2018joint} & Cine sequence & Unsupervised motion branch; weakly supervised segmentation & Shared multi-scale encoder; Siamese recurrent spatial transformer & Motion field and segmentation\\
Apparent flow & 2019, Med Image Anal~\citep{zheng2019apparent} & 2D$+t$ cine, two phases & Semi-supervised: many unsegmented plus few expert-segmented images & Pixel-wise apparent flow combined with segmentation masks & Flow maps; per-segment radius and thickness time series; pathology class\\
Motion Pyramid Networks & 2020, MICCAI~\citep{yu2020mpn} & Cine image pair & Unsupervised, with cyclic teacher--student distillation & Multi-scale pyramid of motion fields, fused & Refined motion field in a single step\\
FOAL & 2020, CVPR~\citep{yu2020foal} & Cine sequence at test time & Meta-learned online optimiser & Test-time adaptation of a learned tracker to the new sequence & Adapted motion field (about 0.4~s per video)\\
MulViMotion & 2022, IEEE TMI~\citep{meng2022mulvimotion} & Short-axis plus two- and four-chamber long-axis cine & Unsupervised, with weak shape supervision & Hybrid 2D/3D network; shape regularisation module during training & Dense 3D motion field of the LV myocardium\\
DeepMesh & 2024, IEEE TMI~\citep{meng2024deepmesh} & Short- and long-axis cine & Unsupervised, via a differentiable mesh-to-image rasteriser & Template 3D mesh of endocardial and epicardial surfaces propagated to subject space & Vertex-wise 3D displacement with maintained vertex correspondence\\
Low-rank groupwise deformations & 2024, preprint (MSc thesis)~\citep{rendell2024lowrank} & Image sequence and a target image & Unsupervised & Groupwise diffeomorphic registration driven towards a low-rank registered stack & Groupwise deformations\\
MCM (Mamba) & 2025, MICCAI RIME workshop~\citep{yin2025mcm} & Target image sequence from the cardiac cycle & Unsupervised & Bi-directional Mamba block with bi-directional scanning; motion decoder using adjacent frames & Temporally consistent deformation fields\\
Cardio\-Morph\-Net & 2026, Med Image Anal~\citep{movahed2026cardiomorphnet} & Short-axis cine volumes & Recursive registration of segmentation maps; no intensity similarity loss & Recurrent variational autoencoder with separate posteriors for segmentation and motion & 3D motion field plus a voxel-wise uncertainty map\\
Bidirectional latent neural ODE & 2026, Comput Med Imaging Graph~\citep{ruan2026latentode} & Cine sequence & Unsupervised & Continuous-time latent ODE with bidirectional forward--backward evolution and Lagrangian consistency & Inter-frame deformation fields over the cycle\\
STCMT-Net & 2026, Med Image Anal~\citep{qiao2026stcmt} & 4D cardiac images (computed tomography and cine MR) & Unsupervised & Keypoints detected without supervision as structural anchors; the dense field is a weighted combination of basis vectors from matched keypoint pairs, with residual refinement & Dense motion field; derived strain\\
\midrule
\multicolumn{6}{@{}l}{\textit{Physics-, biomechanics- and representation-constrained estimators}}\\
Biomechanics-informed prior & 2023, Med Image Anal~\citep{qin2023biomech} & Cine sequence & Unsupervised; prior learned from simulated deformations & Variational autoencoder manifold of biomechanically plausible deformations, traversed at tracking time & Motion field and derived strain\\
WarpPINN & 2023, Med Image Anal~\citep{arratia2023warppinn} & Cine bSSFP, fitted per case & None (per-case optimisation) & Neural network deformation with hyperelastic regulariser and Jacobian-based near-incompressibility; Fourier feature mapping against spectral bias & Deformation field; strain by automatic differentiation\\
WarpPINN-fibers & 2026, Eng Comput~\citep{alvarez2026warppinn} & Cine MR & None (per-case optimisation) & As WarpPINN plus a fibre-stretch penalisation along synthetically generated fibres & Deformation field and fibre-informed strain\\
Implicit neural representation & 2026, STACOM~\citep{bell2026inr} & not reported & not reported & Implicit neural representation of intramyocardial motion and strain & not reported\\
Space-time finite-element correlation & 2026, Med Image Anal~\citep{liu2026fe} & Routine cine, short- and long-axis & None (correlation framework) & Region-specific biomechanical regularisation, data-driven temporal decomposition, correlation-based multi-view alignment & 2D/3D$+t$ motion and strain\\
\midrule
\multicolumn{6}{@{}l}{\textit{Reference-supervised and end-to-end strain}}\\
DeepTag & 2021, CVPR~\citep{ye2021deeptag} & Tagged MR (comparator acquisition) & Unsupervised & Bidirectional diffeomorphic registration between consecutive frames; differentiable composition to a Lagrangian field & Inter-frame and Lagrangian motion fields\\
StrainNet & 2023, Radiol Cardiothorac Imaging~\citep{wang2023strainnet} & Time series of myocardial contours (from DENSE magnitude at training, cine feature tracking at test) & Supervised on DENSE displacement & 2D$+t$ convolutional network; prior is implicit in the DENSE-derived training distribution & Intramyocardial displacement; global and segmental circumferential strain\\
DENSE-guided rotation network & 2025, Proc IEEE ISBI~\citep{lei2025denserotation} & Routine cine bSSFP & Supervised on DENSE displacement & Separate time-series rotation branch and deformable-registration radial branch, combined in a fusion network & Myocardial displacement, with large rotation modelled explicitly\\
MyoNet & 2026, Bioengineering~\citep{an2026myonet} & Cine (pre-clinical rat model) & Supervised against a sine-wave-modelling reference & Convolutional network with alternating kernel sizes for spatial and temporal analysis & Displacement maps; circumferential and radial strain\\
DeepStrain & 2021, Front Cardiovasc Med~\citep{morales2021} & Cine bSSFP & Image, anatomical-label and smoothness losses & CarSON segmentation plus CarMEN motion network; end-diastolic reference & Segmentation, displacement, Green--Lagrange radial and circumferential strain\\
Deep learning synthetic strain & 2023, Radiol Cardiothorac Imaging~\citep{masutani2023dlss} & Cine & Supervised on a radiologist wall-motion reference & not reported in detail & Regional wall-motion assessment from synthetic strain\\
Deep learning aligned strain & 2025, Radiol Artif Intell~\citep{koehler2025alignedstrain} & Cine & Supervised; LGE-defined fibrosis as the reference standard & not reported in detail & Aligned strain used to flag fibrotic myocardium\\
Cine-only infarct delineation & 2018, Med Image Anal~\citep{xu2018} & Non-contrast cine sequence & Supervised on manual delineation of delayed-enhancement images & Long short-term memory motion features plus optical flow, fused per pixel & Pixel-wise infarct map\\
Survival-supervised motion & 2019, Nat Mach Intell~\citep{bello2019} & Cine-derived motion & Supervised on survival outcome & Learned low-dimensional motion representation & Risk prediction\\
\bottomrule
\end{longtable}

\begin{longtable}{@{}L{0.145\textwidth}L{0.245\textwidth}L{0.275\textwidth}L{0.282\textwidth}@{}}
\caption{How each estimator has been validated, in the row order of \cref{tab:methods_what}. The
last column is the one that bears directly on this review's question. \texttt{not reported} means
the cited source does not state the item; no entry has been inferred from a related
publication.}\label{tab:methods_validation}\\
\toprule
Method & Training and evaluation data & Validation approach reported & Evaluation on pathology or focal abnormality\\
\midrule
\endfirsthead
\toprule
Method & Training and evaluation data & Validation approach reported & Evaluation on pathology or focal abnormality\\
\midrule
\endhead
\multicolumn{4}{@{}l}{\textit{Conventional estimators}}\\
Boundary feature tracking & Many clinical cohorts; the cross-software comparisons are collected in \cref{sec:crosssoftware} & Inter-software and inter-version agreement, scan--rescan precision, comparison with tagging or DENSE, clinical association & Studied in many disease cohorts; regional and radial agreement is consistently weaker than global agreement~\citep{militaru2021,brandt2024,bell2025strain8}\\
Automated segmentation plus spline tracking & 42 patients with hypertrophic cardiomyopathy and 21 healthy controls~\citep{vigneault2017} & Segmentation accuracy against expert manual segmentation (pixel accuracy about 97\%); endocardial circumferential strain compared between patients and controls & Group separation between hypertrophic cardiomyopathy and controls; no motion reference and no focal abnormality experiment\\
Free-form deformation (FFD) & Three-dimensional contrast-enhanced breast MRI in volunteers and patients & Comparison with rigid and affine registration for recovery of breast motion and deformation & Not applicable; the original study is outside cardiac imaging. Included as the transformation model reused by later cardiac work\\
Groupwise $n$D$+t$ B-splines & Synthetic data and 3D$+t$ lung computed tomography; 3D$+t$ cardiac computed tomography angiography; demonstrated on 2D$+t$ ultrasound and MR & Sub-voxel registration accuracy on synthetic data; LV volume curves deviating on average about 3\% from manual annotations & not reported\\
Groupwise locally low-rank & Simulation ($n=20$), public cine ($n=100$), in-house tagging-MRI patients ($n=16$) & Point-tracking, voxelwise strain and global strain error against simulation; contour-tracking error and late-diastolic drift on public cine; correlation with tagging radial strain in patients & Patient dataset with paired tagging; disease composition not reported in the record retrieved. No prescribed focal abnormality\\
Anatomy-Aware Tracker & Long-axis cardiac cine & Tracking performance against other methods, with qualitative assessment of shape realism & not reported\\
Laplace-based murine framework & Mouse cine MRI & Comparison of derived strain with tagged MR in the same animals & Not a human cine study; focal abnormality not reported\\
Cleared nearly incompressible model & Non-clinical verification using analytical and realistic phantoms, real MR data and expert manual tracking; the record states that no clinical study was needed for substantial equivalence & Regulatory substantial-equivalence evidence~\citep{fdaK232661} & Not established by the record; a 510(k) summary states device identity and intended use, not regional accuracy\\
Harmonic phase (HARP) & Real and simulated tagged MR images & Synthesis of conventional tag lines; displacement reconstruction for small motions; 2D strain & Not applicable; tagged acquisition, included as the material-sensitive comparator\\
\midrule
\multicolumn{4}{@{}l}{\textit{Unsupervised learning-based registration}}\\
VoxelMorph & General medical image registration datasets & Comparison of registration accuracy and run time against the methods the authors selected as comparators; effect of training-set size & not reported\\
TransMorph & Inter-patient and atlas-to-patient brain MRI; phantom-to-computed-tomography & Comparison with existing registration methods and Transformer architectures; topology preservation in the diffeomorphic variant; calibration of the Bayesian uncertainty & not reported for cardiac pathology\\
Fourier-Net & Image-registration benchmarks & Registration accuracy with inference speed, memory and multiply-add operations & not reported\\
Fourier-Net+ & Three registration datasets & Accuracy comparable with the authors' selected comparator methods at lower computational cost & not reported\\
\midrule
\multicolumn{4}{@{}l}{\textit{Cardiac-specific estimators}}\\
Joint motion and segmentation & Cardiac MRI from 220 subjects & Accuracy and speed against competing methods for both tasks & not reported\\
Apparent flow & ACDC & Classification accuracy of 95\% (training) and 94\% (testing) from nine motion and shape features & Pathology classes are the target; the motion field itself is validated through the downstream label, not against a motion reference\\
Motion Pyramid Networks & Two public clinical datasets & Multiple metrics plus inference time against a strong baseline; additional clinically oriented error metrics proposed & not reported\\
FOAL & Two public clinical datasets & Accuracy against an offline-trained tracker under train--test distribution mismatch; 0.4~s per video & Distribution shift is the explicit target, but the shift is not characterised as focal myocardial abnormality\\
MulViMotion & 580 UK Biobank subjects & Quantitative and qualitative comparison with competing methods for 3D motion tracking of the LV myocardium & not reported\\
DeepMesh & UK Biobank & Mesh-based 3D motion estimation of the LV, with vertex correspondence maintained across frames & not reported\\
Low-rank groupwise deformations & not reported & Comparison of resulting low-rankness with other groupwise approaches & not reported\\
MCM (Mamba) & Two public datasets & Quantitative and qualitative comparison with conventional and learning-based trackers & not reported\\
Cardio\-Morph\-Net & UK Biobank and M\&Ms & Warped mask shapes compared with ground-truth masks; uncertainty values compared with other probabilistic methods & M\&Ms is multi-centre and multi-pathology, but the published comparison is mask agreement, not focal motion accuracy\\
Bidirectional latent neural ODE & ACDC and M\&Ms & Comparison with the authors' selected comparator methods for temporal consistency and plausibility of the motion field & not reported\\
STCMT-Net & 4D computed tomography and cine MRI datasets & Comparison with competing motion-estimation methods, plus a strain analysis reported to separate patients under different cardiac conditions & Group separation by cardiac condition is reported; no prescribed focal abnormality and no motion reference\\
\midrule
\multicolumn{4}{@{}l}{\textit{Physics-, biomechanics- and representation-constrained estimators}}\\
Biomechanics-informed prior & Three public cardiac cine datasets & Motion-tracking accuracy, volume preservation, generalisability across data distributions, and derived strain & not reported as a prescribed focal abnormality experiment\\
WarpPINN & Synthetic examples and a cine bSSFP benchmark of 15 healthy volunteers & Landmark tracking against existing methods; strain obtained by automatic differentiation & Healthy benchmark; focal abnormality not evaluated\\
WarpPINN-fibers & Synthetic phantom and the same 15-participant healthy benchmark & Strain estimation with fibre information compared with the fibre-free formulation & Healthy benchmark; focal abnormality not evaluated\\
Implicit neural representation & not reported & not reported & not reported\\
Space-time finite-element correlation & One synthetic dataset with ground-truth motion and strain; three public datasets with ground-truth landmarks or masks; one clinical dataset with ground-truth masks & 2D$+t$ motion and strain accuracy and temporal consistency against two classical feature-tracking and four deep-learning baselines; multi-view alignment and 3D$+t$ motion on the clinical dataset only & One synthetic ground-truth setting; the in-vivo references are landmarks and masks. Systematic variation of abnormality size or severity is not reported\\
\midrule
\multicolumn{4}{@{}l}{\textit{Reference-supervised and end-to-end strain}}\\
DeepTag & A representative clinical tagging dataset & Landmark tracking accuracy and inference efficiency against conventional tracking & not reported\\
StrainNet & 161 patients, 99 healthy adults, 45 children and adolescents, with paired DENSE & End-point error of $0.75\pm0.35$~mm against DENSE; ICC with DENSE 0.87 (global) and 0.75 (segmental) circumferential strain, versus 0.72 and 0.48 for commercial feature tracking & Segmental agreement against a material-sensitive reference in a mixed patient population; abnormality size and severity were not varied experimentally\\
DENSE-guided rotation network & not reported in the record retrieved & Strain accuracy against DENSE ground truth, with large rotation as the stated target & not reported\\
MyoNet & Dahl salt-sensitive rat models undergoing radiation therapy & Agreement with strain derived by sine-wave modelling: ICC 0.973 (circumferential) and 0.975 (radial), with structural-similarity and correlation measures & Radiation-therapy injury in a rodent model; not human cine, and the reference is analysis-derived rather than prescribed\\
DeepStrain & ACDC ($n=150$, including the 50-subject official segmentation test set with predicted labels); CMAC; ten healthy volunteers & CMAC landmark end-point error $2.9\pm1.5$~mm; scan--rescan repeatability; later independent synthetic evaluation (see \cref{sec:focal}) & Independently evaluated on MRXCAT2.0 synthetic normal, dilated, hypertrophic and infarcted cases; see \cref{sec:focal} for the reported errors\\
Deep learning synthetic strain & 53 patients, 846 segments & AUC 0.90 against the consensus of four cardiothoracic radiologists; sensitivity, specificity and accuracy 86\%, 85\% and 86\% at a 30\% peak radial strain threshold; reader kappa 0.60--0.78 & Regional wall-motion abnormality is the target, but the reference is expert reading, not measured motion\\
Deep learning aligned strain & 139 male patients with Duchenne muscular dystrophy; fibrosis evaluation in 57, reproducibility in 82 & Specificity $+40\%$, accuracy $+17\%$, and accuracy in patients with preserved ejection fraction $+61\%$ relative to the comparator & Target is LGE-defined fibrosis, a tissue reference rather than a motion reference\\
Cine-only infarct delineation & 165 cine datasets & Pixel accuracy 95.03\%, Dice 89.87\% and $\kappa=0.91$ against manual delineation of delayed-enhancement images by two radiologists & Infarction is the target, but the reference is LGE-defined tissue, not motion\\
Survival-supervised motion & 302 patients & Harrell C index 0.75 (95\% CI 0.70--0.79) versus a human benchmark of 0.59 (0.53--0.65), $P=0.0012$ & Prognostic discrimination; neither regional motion accuracy nor focal abnormality was evaluated\\
\bottomrule
\end{longtable}

Three patterns in \cref{tab:methods_validation} are worth stating explicitly, because they define the
rest of this review. First, the modal evidence across the methods reviewed is agreement between a
propagated and a reference segmentation, supplemented by sparse landmark error; this is reasonable
evidence for the claims those papers make and, by the construction in \cref{fig:correspondence}, it
cannot by itself certify material correspondence. Second, where a method in that table is evaluated
against a material-sensitive reference, the reference is almost always acquired in a mixed or
healthy population, and the comparison is reported as agreement rather than as error under controlled
conditions. Third, the right-hand column is mostly \texttt{not reported}: among the estimators
reviewed, deliberate evaluation on a focal abnormality of known size and severity appears in one
independent synthetic study, and the evidence that does exist on pathological hearts tends to be
mask agreement or a downstream label. This is a statement about the published evaluation record, not
about the methods' capability, and it is exactly the gap that \cref{sec:regional,sec:synthesis}
characterise.

\section{How motion and strain estimates have been validated}\label{sec:validation}

\subsection{What each validation target can establish}\label{sec:targets}

The previous section compared estimators by what they do. This section classifies the evidence by
what it can support. A study may test agreement with observed anatomy, consistency between
estimators, error against prescribed motion, agreement with a material-sensitive acquisition,
repeatability, association with tissue, or association with outcome. These targets are not
interchangeable, and a result at one point in the measurement chain does not retrospectively
establish the steps before it. \Cref{tab:targets} separates, for each target, the metric or record it
produces, the interpretation it supports and the questions it leaves open. \Cref{fig:targets} carries
the same classification as a single display, setting what each target \emph{can} support directly
against what it does not establish on its own; that two-column contrast is the order in which
\cref{sec:overlap} to \cref{sec:unit} then work through the evidence.

\begin{longtable}{@{}L{0.16\textwidth}L{0.20\textwidth}L{0.26\textwidth}L{0.29\textwidth}@{}}
\caption{Validation targets and what they support. The rows are complementary questions, not an
accuracy ranking. Regulatory records are listed as a distinct source of device claims, not as a
motion reference.}\label{tab:targets}\\
\toprule
Target or reference & Typical metric or record & Supported interpretation & Questions it leaves open\\
\midrule
\endfirsthead
\toprule
Target or reference & Typical metric or record & Supported interpretation & Questions it leaves open\\
\midrule
\endhead
Anatomical overlap & Dice, surface distance, contour error & Propagated anatomy aligns with a reference segmentation & Tangential material correspondence, dense interior motion and the accuracy of strain gradients\\
Image-registration similarity & Image-similarity score or registration data term & Agreement with image intensities under the chosen matching criterion & Material trajectories; independent accuracy, when the same criterion is both optimised and evaluated\\
Inter-estimator or inter-software agreement & Bias, limits of agreement, intraclass correlation & Consistency of specified outputs between the tested implementations & Accuracy against material motion, whether a shared bias is present, and equivalence under another strain definition\\
Sparse landmarks & Endpoint or trajectory error with observer uncertainty & Correspondence near the selected visible or tagged points & Error between landmarks, and derivative accuracy through the wall\\
Implanted physical markers & Segment shortening or crystal-pair distance against the imaging estimate & Agreement with a physical material reference inside a living heart under the studied preparation & Human applicability, sparse spatial sampling, and the effect of the surgical preparation on mechanics\\
Prescribed motion & Vector error, strain bias, localisation, width, spillover, timing, topology & Technical accuracy under the stated simulated geometry, motion and image formation & In-vivo accuracy and behaviour outside the simulated regime\\
Tagging or DENSE & Displacement or strain agreement after spatial and temporal registration & Agreement with an independent material-sensitive acquisition under the paired conditions & Error-free truth, performance at untested sites, and outcome benefit\\
Scan--rescan & Within-subject standard deviation, coefficient of variation, repeatability coefficient & Precision under the specified scanner, protocol, software, reader and interval & Accuracy, absence of systematic bias, and transfer to another version or population\\
Independent tissue target (LGE) & Zone contrast, discrimination, spatial association & Relationship between a strain marker and independently defined tissue regions & Dense motion truth, a causal mechanism, and one-to-one tissue--strain correspondence\\
Clinical outcomes & Adjusted association, discrimination, calibration, decision analysis & Prognostic or diagnostic association, or decision-related performance as directly studied & Material-motion accuracy, threshold transport, and treatment benefit without the corresponding trial design\\
Regulatory record & Intended use, technical comparison, submitted verification & Device status and bounded claims within one jurisdiction & Superiority, outcome benefit, interchangeability and routine adoption\\
\bottomrule
\end{longtable}

\begin{figure}[!ht]
 \centering
 \includegraphics[width=\textwidth,keepaspectratio]{figures/p3_Figure3_validation_targets.pdf}
 \caption{Validation targets support different claims. Each row names a reference or comparison, what
 it can support, and what it does not establish on its own. The rows are complementary rather than
 ranked, and none describes the performance of any particular estimator.
 \enquote{Independently defined} late gadolinium enhancement (LGE) regions means that the tissue
 regions are defined separately from the strain marker; it does not assert statistical independence
 or a dense motion truth. The implanted-physical-marker row of \cref{tab:targets} has no counterpart
 in this figure; it is discussed in \cref{sec:focal}.}
 \label{fig:targets}
\end{figure}

\subsection{Anatomical and image consistency}\label{sec:overlap}

Anatomical overlap tests whether a propagated structure matches a reference segmentation, and the
annulus construction of \cref{sec:illposed} shows the limit of that test exactly: a smooth
one-to-one reparameterisation can leave Dice at 1 and the Hausdorff distance at 0 while the
tangential stretch varies by a factor of three. The counterexample is analytic and says nothing
about how often any particular estimator behaves this way; it establishes only that overlap cannot
be used as a substitute. The same applies to the data term of \cref{eq:registration}: a lower
registration loss is evidence of a better fit to that criterion, not an independent measurement of
trajectory accuracy, and where the myocardial signal does not distinguish between candidate
transformations, similar data-term values can accompany different displacement gradients.

\subsection{Prescribed motion}\label{sec:prescribed}

In a design that uses images alone, a prescribed-motion experiment is the setting in which regional
error can be measured directly, because the deformation is known by construction; the one other
material reference, physical markers implanted in the tissue, requires surgery and is taken up in
\cref{sec:focal}. Several generators now exist, and they make
different trade-offs between realism and control. Adams et al.\ built a single-slice numerical
phantom from one volunteer's acquisitions with the JEMRIS open-source MR
simulator, at in-plane resolutions from
$1.4\times1.4$ to $3.0\times3.0$~mm$^2$ and 20 to 40 cardiac phases, and used it to test two
feature-tracking packages against known radial and circumferential strain. Increasing pixel size
produced a trend of increasing departure from the ground-truth peak strain, reaching 9.17 percentage
points for radial and 8.42 percentage points for circumferential strain for one cvi42 version
between the finest and coarsest resolutions~\citep{adams2023}. Mukherjee et al.\ derived an
in-silico phantom from finite-element simulations, first recovering an exact twist angle for a
thick-walled cylinder under pure torsion and then synthesising dynamic images from mouse-specific
simulations~\citep{mukherjee2024}.

Two of the generators carry pathology. MRXCAT2.0 couples the XCAT torso phantom with a statistical shape
model and a biophysical electromechanical model, generating paired bSSFP images and ground-truth LV
function across healthy, infarcted, dilated and hypertrophic cases~\citep{mrxcat2023}. The Multimodality
STRAUS framework generates 3D synthetic cine MR, tagged MR and ultrasound sequences for the same
virtual patient from an electromechanical model, and its released database is the most
pathology-dense prescribed-motion resource in \cref{tab:datasets}: three anatomical templates, each
simulated under six conditions, give 18 virtual patients of which 12 are ischaemic across four coronary
configurations, with 90 sequences and per-modality reference meshes released publicly and without
login~\citep{zhou2018straus}. Two properties bound its use. The images are synthesised from the
prescribed deformation, so an estimator whose own deformation prior resembles the generator's can be
flattered by construction, and the released set is a fixed list of cases rather than a parameterised
generator, so scar location and size cannot be swept without re-running the underlying
electromechanical pipeline.

DENSE-SIM targets the reference acquisition rather than cine. It generates cine DENSE images with
high-resolution known strain, and its released example of 180 short-axis slices was used to
characterise how the spatial regularisation parameter of a standard DENSE analysis tool trades bias
against variance. At a typical setting the average signed error in end-systolic Green radial strain
was $0.04\pm0.09$ and circumferential $-0.00\pm0.04$, with radial strain showing the larger
underestimation, particularly near the endocardium~\citep{barbaroux2025,densesimrepo}.

The DENSE-SIM result deserves emphasis because it is about the reference itself. Even a
material-sensitive acquisition has a regularisation parameter whose setting biases regional radial
strain, so the uncertainty of the reference belongs in any comparison that uses it. Simulator
evidence as a whole is bounded by its anatomy, motion, texture and acquisition assumptions, and
visual realism does not by itself establish in-vivo validity. \Cref{tab:datasets} lists the public
resources that a regional cine validation study could currently draw on.

\begin{longtable}{@{}L{0.17\textwidth}L{0.19\textwidth}L{0.20\textwidth}L{0.17\textwidth}L{0.18\textwidth}@{}}
\caption{Public datasets and simulation frameworks relevant to regional cine motion validation.
\enquote{Motion reference} means a reference for material displacement or strain, not a
segmentation. Entries describe the resource as documented in the cited record.}\label{tab:datasets}\\
\toprule
Resource & Modality and content & Population & Motion reference & Pathology represented\\
\midrule
\endfirsthead
\toprule
Resource & Modality and content & Population & Motion reference & Pathology represented\\
\midrule
\endhead
CMAC~\citep{cmac2011,tobongomez2013}; used by DeepStrain in \cref{tab:methods_validation}~\citep{morales2021} & 3D tagged MR, cine bSSFP and 3D ultrasound (1,158 image volumes) & A dynamic phantom and 15 healthy volunteers & 12 landmarks (four walls at three ventricular levels) tracked over the cycle in the 3D tagged data by two observers; median inter-observer variability 0.77~mm (phantom) and 0.84~mm (volunteers) & None; healthy volunteers only\\
Stanford \emph{Human MRI Data to Quantify Cardiac Performance}~\citep{stanforddata} & Cine bSSFP ($1\times1\times8$~mm, 25 phases, contiguous short-axis stack covering the left ventricle plus five long-axis views rotated $37^\circ$ apart), xyz-encoded cine DENSE \emph{at slice locations identical to the cine}, three mid-ventricular short-axis grid-tagging slices (30 phases, 7~mm tag spacing) and native T1 at those same locations, plus cardiac diffusion tensor imaging; de-identified DICOM, one archive per subject; Open Data Commons Open Database Licence v1.0 & 55 participants imaged 2023--2025 on one 3-T system with identical coils and protocol: 51 healthy volunteers (mean age $39\pm14$ years) and 4 patients with aortic stenosis & DENSE displacement at corresponding slice locations; a processed measurement, not a prescribed truth. Cine is retrospectively gated and DENSE prospectively gated, so the two must be aligned in time & No focal infarction. The only pathology is aortic stenosis in 4 of 55, which is a pressure-overload rather than a regional-injury phenotype\\
Multimodality STRAUS~\citep{zhou2018straus} & Synthetic 3D cine MR, tagged MR in three channels and 3D ultrasound for every virtual patient; 90 sequences and 2,700 image volumes as \texttt{raw}/\texttt{mhd} with per-modality \texttt{vtk} reference meshes; public with no login, and no licence stated on the project page & 18 virtual patients: three anatomical templates, each simulated under six conditions (healthy, left bundle branch block, and ischaemia in the left circumflex, left anterior descending and two right coronary configurations), giving 3 healthy, 3 dyssynchronous and \textbf{12 ischaemic} cases & Motion fields prescribed by an electromechanical model, supplied as reference meshes; the images are then generated from those fields, so an estimator whose deformation prior resembles the generator's carries an inverse-crime risk & Ventricular dyssynchrony and myocardial ischaemia across four coronary territories. The released set is a fixed list of cases: infarct location and size are not exposed as adjustable parameters\\
MRXCAT2.0~\citep{mrxcat2023} & Synthetic bSSFP cine with paired ground truth & Shape model trained on a clinical population; four demonstrated use cases & Prescribed biophysical deformation and strain & Normal, infarcted (one elliptical free-wall scar), dilated and hypertrophic\\
DENSE-SIM~\citep{barbaroux2025,densesimrepo} & Synthetic cine DENSE with sub-voxel ground-truth strain & 180 short-axis slices in the released example & Prescribed Lagrangian strain & Not stated; the released example varies motion and signal-to-noise ratio\\
Numerical FT phantom~\citep{adams2023} & Single-slice simulated cine at several resolutions & Derived from one volunteer & Prescribed radial and circumferential strain & None; resolution is the varied factor\\
In-silico heart phantom~\citep{mukherjee2024} & Synthetic MR from finite-element simulations & Cylinder plus mouse-specific models & Exact analytic twist; simulated kinematics & Not reported for human pathology\\
ACDC~\citep{bernard2018acdc} & Cine with expert segmentations, reference measurements and an expert diagnostic label & 150 multi-equipment examinations, 100 training and 50 test & None & Five equally sized classes: previous myocardial infarction, dilated and hypertrophic cardiomyopathy, abnormal right ventricle and normal\\
UK Biobank CMR~\citep{petersen2016ukb} & Cine (short- and long-axis) in a population cohort & Tens of thousands of participants & None & Population-based; disease is incidental\\
Emidec~\citep{lalande2020emidec} & Delayed-enhancement MR only, in the short axis, with expert myocardial and diseased-area contours plus emergency-department clinical parameters & Number of examinations not stated in the record retrieved & None; this is a tissue resource, not a motion resource & Normal cases and patients with myocardial infarction\\
Cardiac Atlas Project~\citep{fonseca2011cap,suinesiaputra2018} & Imaging examinations with associated clinical data and 3D LV surface-point shape models & Multi-cohort; the infarct-classification challenge used 100 patients with myocardial infarction and 100 asymptomatic volunteers for training and 200 further cases for validation & None & Myocardial infarction versus asymptomatic\\
CMRxRecon2024 \citep{wang2025cmrxrecon2024} & Multi-contrast, multi-view raw $k$-space including cine and tagging & Not read in this review; see note & None (reconstruction benchmark) & Not read in this review\\
\bottomrule
\end{longtable}

\noindent\emph{Notes to \cref{tab:datasets}}: the CMRxRecon2024 record was verified bibliographically
but its abstract text could not be retrieved in this review, so its participant numbers and split
sizes are deliberately not stated here. The STRAUS and Stanford entries were re-checked on
8~October 2026, the first against the project page and the database catalogue, the second against
the Stanford Digital Repository record and the laboratory protocol page. The case
composition of STRAUS is taken from the catalogue listing rather than inferred, and the Stanford
cohort is given as the repository record states it, which includes four participants the laboratory
protocol page does not mention.

\subsection{Tagging and DENSE as in-vivo references}\label{sec:material}

Tagging and DENSE supply material-sensitive information in vivo, and comparisons against them are
the closest thing to a material reference that is routinely available in human studies. They remain
conditional on the acquired components, on spatial and temporal registration between the two
acquisitions, and on the strain definition used for each. Wehner et al.\ analysed 88 participants and 186 paired slices and found
that some feature-tracking whole-slice strain measures agreed with DENSE and with simple two-contour
geometry, while basal, apical, torsional and dyssynchrony measures were less
reliable~\citep{wehner2018}; agreement of a whole-slice contour-derived measure can occur without
establishing that interior points were followed correctly. Cao et al.\ reported systematic
differences among DENSE, tagging and feature tracking~\citep{cao2018}, which is a reminder that
dedicated motion-sensitive acquisitions are neither interchangeable with one another nor error-free.

How fallible the reference is has itself been measured, and the answer is encouraging for the
reference and relevant to every comparison that uses it. Li et al.\ repeated breath-hold
two-dimensional cine DENSE in a mid-ventricular short-axis slice before and after contrast in two
cohorts of 11, one receiving gadolinium and one receiving three incremental doses of ferumoxytol,
and assessed circumferential strain against a declared region of practical equivalence of $\pm0.02$
in a Bayesian analysis. Global circumferential strain showed practical equivalence in both cohorts,
and segmental reproducibility was high throughout, with at least 87.02\% of comparisons falling
inside that region~\citep{li2025densereprod}. Two readings follow, and they point in opposite
directions. The reference is more stable under a deliberate perturbation than the phrase
\enquote{not an error-free truth} on its own suggests, which strengthens rather than weakens
comparisons made against it. At the same time the study measures reproducibility across a contrast
manipulation in 22 subjects at one site, not accuracy, and the regularisation result of
\cref{sec:prescribed} shows that a processing choice inside the same pipeline still moves regional
radial strain. Reference stability and reference accuracy remain separate properties.

Resources that pair the two acquisitions in the same participants are what make such comparisons
possible at all, and the most completely specified one currently public is the Stanford
\emph{Human MRI Data to Quantify Cardiac Performance} collection. Its value for this review's
question lies in a single acquisition decision: the cine DENSE slices, short-axis and long-axis
alike, were prescribed at locations identical to the cine slices, and the three mid-ventricular
tagging slices and native T1 maps sit at those same locations, so a cine estimate and a
material-sensitive reference can be compared without a between-slice interpolation step. The
collection holds 55 participants imaged on one 3-T system between 2023 and 2025 with identical coils
and protocol, of whom 51 are healthy volunteers and 4 have aortic stenosis, and it is released as
de-identified DICOM under an Open Data Commons Open Database Licence~\citep{stanforddata}; the
custodians recommend a published tag-tracking pipeline for the tagging
series~\citep{loecher2021tagging}. Three properties bound what it can establish. It contains no
focal infarction, so it cannot show whether a regional deficit survives estimation. Its cine is
retrospectively gated while its DENSE is prospectively gated, so the two must be aligned in time and
late diastole may be missing from one of them. And DENSE displacement is itself the output of phase
unwrapping, contouring and spatial regularisation, so a comparison against it supports a statement
about \emph{agreement with DENSE} and not about accuracy.

Militaru et al.\ provide the most directly scale-resolved comparison. In 61 participants under one
3-T protocol, agreement with tagging gave global longitudinal-strain ICCs of 0.92--0.94 and global
circumferential-strain ICCs of 0.88--0.91, whereas the corresponding regional ranges were 0.44--0.50
and 0.57--0.77, and regional radial-strain ICCs were 0.05--0.51~\citep[Detailed Results and
Fig.~5]{militaru2021}. The study describes a two-way mixed-effects model for overall agreement; a
single-measure versus average-measure designation was not identified in the available methods, and
the detailed-results ranges are retained here because the abstract gives partly different ranges
without an identified explanation for every discrepancy. No adequacy threshold is applied to these
coefficients here, because \cref{sec:recoverability} declines to set a tolerance until the decision a
measurement serves has been named; what the numbers carry is the within-study contrast between the
global and the regional value, which is internal to one cohort and one protocol and is therefore
robust to the dependence of an intraclass correlation on the between-segment variance of the sample.
Brandt et al.\ reached a compatible
conclusion in 42 subjects comparing MRI tagging, MRI feature tracking and ultrasound speckle
tracking: global peak strain ICCs of 0.71--0.83 without significant bias, but significant biases in
all segmental peak-strain parameters, with segmental ICCs of 0.62--0.77 except for lateral
longitudinal strain at 0.23 ($P=0.22$)~\citep{brandt2024}.

Kihlberg et al.\ compared the acquisitions against an independent tissue target rather than against
each other: in 116 patients with suspected coronary artery disease, for detecting segments with more
than 50\% LGE transmurality at 80\% specificity, sensitivity was 82\% for DENSE and 71\% for
tagging, and segment-level AUCs were 0.87 for DENSE, 0.83 for tagging ($P<0.001$) and 0.66 for
feature tracking ($P=0.003$)~\citep{kihlberg2020}. This ranks the acquisitions by their ability to
identify tissue-defined dysfunction, which is a clinically meaningful ordering, and it is not a
measurement of motion accuracy for any of the three.

The post-cutoff study of Jahromi et al.\ illustrates how much the comparison depends on the
measurement definition. In 40 healthy volunteers contributing two short-axis slices each, alongside
an idealised cylindrical phantom, cine showed closer endocardial circumferential agreement,
lower-magnitude epicardial circumferential strain, and higher-magnitude radial and longitudinal
strain than two-slice DENSE; the authors reported equivalence for peak slice-wise endocardial
circumferential strain within a selected margin of $\pm3.5$ percentage points, with a mean
cine-minus-DENSE difference of $-2.7$ percentage points~\citep{jahromi2026densecine}. This is a
selected-margin comparison of slice-wise peaks in healthy hearts, with two dependent slices per
participant; it does not address individual interchangeability or accuracy within an infarct region.

\subsection{Repeatability}\label{sec:repeatability}

Repeatability addresses precision under repeated measurement conditions, which is a different
property from accuracy. The Quantitative Imaging Biomarkers Alliance framework separates bias,
repeatability and reproducibility and ties each to a defined measurand and use
condition~\citep{raunig2015}. In the same-day strain-8 experiment, 20 healthy volunteers underwent
eight CMR and echocardiographic protocols; repeatability varied by direction and by spatial scale,
with many segmental measures less precise than global ones, and for 1.5-T CMR feature-tracking
global longitudinal strain the reported coefficient of variation was 9.0\% (95\% confidence interval
6.0--11.2\%)~\citep{bell2025strain8}. A confidence limit for a coefficient of variation is not a
threshold for change between two individual measurements, and 11.2\% is not interpreted here as a
validated individual-change threshold. The quantity a clinician needs is the repeatability coefficient, which under the usual assumptions
is $2.77$ times the within-subject standard deviation and bounds the difference expected between two
measurements of an unchanged subject; the strain-8 segment-level values were not extracted at the
reading depth used here and are therefore not reproduced. A repeatability coefficient and a
clinically important change answer different questions: the first asks whether two measurements differ by more than the stated
technical conditions predict, the second whether the difference matters for management, and
biological day-to-day variation adds a third source of change.

\subsection{Consistency across software and versions}\label{sec:crosssoftware}

Agreement between analysis packages addresses the consistency of the tested outputs rather than
their accuracy against an independent reference, and it is nonetheless important here, because a
regional value that changes with the software cannot be compared across studies at all.
\Cref{tab:software} keeps each result tied to its tested versions, population and reference. Across
these studies, global longitudinal and circumferential agreement was generally stronger than
regional or radial agreement, and reproducible measurements were not necessarily
interchangeable~\citep{militaru2021,goetze2026,ochs2026}. Software version matters as much as
vendor: contour propagation, interpolation, coordinate construction, peak detection and default
exclusions can all change between releases while the output label stays the same, which is why
\cref{tab:estimand} asks for the full version alongside the estimand settings. The regulatory record
does not close this gap, because a 510(k) summary establishes device identity, intended use,
predicate and submitted evidence within one jurisdiction, not cross-product
equivalence~\citep{fda2022qib,fdaK232661}.

\begin{longtable}{@{}L{0.20\textwidth}L{0.23\textwidth}L{0.27\textwidth}L{0.24\textwidth}@{}}
\caption{Independent comparison evidence, tied to the versions, population and reference of each
study.}\label{tab:software}\\
\toprule
Study and population & Tested implementations & What the study establishes & Boundary\\
\midrule
\endfirsthead
\toprule
Study and population & Tested implementations & What the study establishes & Boundary\\
\midrule
\endhead
Militaru et al.; 61 participants & cvi42 5.1, Segment 3.0 and TomTec Arena 4.6 & Pairwise global and regional agreement, and comparison with tagging, under one 3-T protocol~\citep{militaru2021} & Global agreement did not transfer to all regional or radial measures; not a current-version comparison\\
Lim et al.; healthy volunteers & Two analysis packages, plus full-stack versus three-slice sampling & Segmentation coverage and software both alter the reported strain under the tested protocol~\citep{lim2021} & Healthy cohort; no outcome threshold and no known-motion regional truth\\
Goetze et al.; 15 travelling healthy volunteers at three 3-T sites & Medis Suite 3.1 and cvi42 5.13 & Scanner and software effects can be separated within the travelling design~\citep{goetze2026} & Small healthy cohort; the detailed table reports centre-specific inter-software global radial strain bias of 18.34--19.93 percentage points, whereas the abstract gives 18.34--19.83\\
Ochs et al.; 183 post-COVID participants and 57 controls & cvi42; CAAS within IntelliSpace Portal 12; fast strain-encoded imaging with MyoStrain 5.2.2 & Method agreement and observer reproducibility~\citep{ochs2026} & No cvi42 version reported; scan--rescan precision was not assessed, and post-COVID evidence does not establish post-infarction accuracy\\
Zhang et al.; heart failure with preserved ejection fraction & Multiple CMR strain packages & Discrimination between patients and controls can differ between software in the studied cohort~\citep{zhang2021hfp} & Group discrimination does not establish prospective prediction, motion accuracy or a treatment rule\\
\bottomrule
\end{longtable}

\subsection{Tissue and outcome association}\label{sec:association}

Two further classes of evidence are frequently cited in support of cine strain and are included here
for completeness, with their interpretation stated. Association with LGE-defined tissue shows that a
strain marker differs between independently defined tissue regions. In 18 participants with chronic
infarction, circumferential-strain discrimination of any LGE had an AUC of 0.802 for HARP tagging
and 0.679 for one feature-tracking package ($P=0.006$), with further significant inter-software
comparisons~\citep[Fig.~7]{militaru2021}; these are segment-level comparisons in a small cohort, with
within-patient dependence and multiple comparisons to consider, and they demonstrate
method-dependent tissue discrimination rather than the accuracy of the underlying motion field.
Mangion et al.\ compared cine feature tracking with DENSE after infarction and found that the
reported incremental prognostic result concerned DENSE circumferential strain rather than the
feature-tracking measure~\citep{mangion2019}, which is a reminder that a prognostic finding belongs
to the acquisition and estimator that produced it. Association with outcome shows prognostic
relevance: the UK Biobank analysis cited in
\cref{sec:introduction} is the largest example, and its analysable denominator differed by strain
component, since segmentation failures prevented global circumferential and radial analysis in 4,390
participants~\citep{chadalavada2024}. Mechanical dispersion derived from echocardiographic
deformation has been associated with ventricular arrhythmia after infarction~\citep{haugaa2010}, and
an inter-vendor echocardiographic study found implementation-dependent dispersion
values~\citep{appadurai2022}; these motivate a marker and an endpoint without supplying a regional
cine motion reference or a transportable CMR threshold.

\subsection{Unit of inference and reference uncertainty}\label{sec:unit}

Two statistical points govern how every number in this section may be read. First, the unit of
inference: segments within a participant are correlated, and analysing them as independent
observations underestimates uncertainty. A high correlation can coexist with a scale offset, and a
non-significant mean difference does not demonstrate equivalence without a declared margin and
sufficient precision, so bias, limits of agreement and the sources of variation are needed alongside
correlation. Repeated scans, readers, scanners and software versions create crossed or nested
variation, and a regional map adds localisation and multiple-comparison
questions~\citep{raunig2015,puntmann2018}. Second, reference uncertainty: tag intersections and
DENSE phase require processing and registration, manual landmarks carry observer variability,
simulated anatomy may not reproduce in-vivo texture, and LGE zones depend on segmentation and
thresholding. The DENSE-SIM regularisation result in \cref{sec:prescribed} quantifies one instance
of this directly. Propagating reference error, or at least performing a sensitivity analysis across
plausible registrations and annotations, shows how much a comparison depends on the reference
itself.

\section{What is known about recovering regional and focal abnormalities}\label{sec:regional}

The previous section classified the evidence by target. This section narrows to the question the
review set out to answer, and asks what is known specifically about a \emph{focal} deficit: its
magnitude, its location, its extent and its timing. Four kinds of study bear on it, and they are
kept apart because they support different claims.

\subsection{Tests against prescribed focal abnormality}\label{sec:focal}

Prescribed-motion experiments and implanted physical markers are the two settings in which regional
error can be measured against a material reference rather than against another estimate: in the first
the deformation is known by construction, and in the second it is measured by a device placed in the
tissue. Implanted markers established the in-vivo credibility of tagging itself. Yeon et al.\ compared
circumferential segment shortening from spatial modulation of magnetisation with sonomicrometry in a
canine coronary-ligation model, with paired data from 19 of 31 studies, and reported that both methods
separated ischaemic from remote myocardium, that tagged shortening matched sonomicrometry in ischaemic
myocardium ($2\pm3$ against $0\pm3$\%, $P=0.067$) and was slightly higher in remote myocardium
($11\pm10$ against $7\pm5$\%, $P=0.033$), with correlations of $r=0.84$ at end-systole and $r=0.88$ in
late systole~\citep{yeon2001sonomicrometry}. That design, a graded regional deficit with a physical
material reference, is the one this review finds missing for cine estimators; it is reported here for a
dedicated tagging acquisition in an animal model, and the surgical preparation, the sparse sampling and
the species all bound its transfer.

Within the cine literature retrieved, the most informative single result is the DeepStrain
evaluation reported in the MRXCAT2.0 paper, which is independent of the DeepStrain developers. Four
use cases were simulated, namely normal, dilated, hypertrophic and infarcted left ventricles, with
ejection fractions of 51\%, 34\%, 41\% and 49\% respectively, and the infarct was modelled as an
elliptical scar in the free wall. In the infarcted simulation, the
ejection fraction stayed at 49\% because remote myocardium compensated for the reduced mobility of
the scar, which is precisely the configuration a global measure cannot see: ground-truth peak
systolic strains were 0.95 (radial), $-0.18$ (longitudinal) and $-0.18$ (circumferential) in remote
tissue against 0.30, $-0.17$ and 0.01 in the scar. The prescribed deficit is therefore strongly
attribute-dependent before any estimator is applied: radial and circumferential strain are largely
abolished inside the scar while longitudinal strain is almost unchanged ($-0.17$ against $-0.18$).
An experiment on these images can test whether an estimator preserves a radial or circumferential
deficit; it cannot test whether a longitudinal deficit survives, because there is almost none to
lose.

Applying DeepStrain to these images gave a mean Dice of 0.82 across cardiac phases, lowest for the
infarcted case and particularly in the thin scar region, and an average displacement error of
$1.0\pm0.9$~mm. For strain, agreement was good for the circumferential component (average error
$0.02\pm0.04$ across all cases) with larger errors in the thinner-walled dilated and infarcted
cases, while radial strain was generally underestimated, reported as $-0.24\pm0.21$ overall with the
lowest performance attributed to the infarcted case at $-0.20\pm0.21$. All values
are averages over three mid-ventricular short-axis slices~\citep{mrxcat2023}.

This result should be read with two qualifications that the source itself supplies. The authors
describe the infarcted case as the worst performer, but the radial error they report for it
($-0.20$) is smaller in magnitude than the across-case average ($-0.24$), and the per-case standard
deviations overlap; the attribution of \enquote{worst} therefore rests on the full set of panels
rather than on those two numbers alone, and should not be quoted as a quantitative ordering. The
paper also states that a simple scar model was implemented and could be extended to arbitrary shapes
and property variations. One scar geometry, one estimator, one slice position and one simulated
signal-to-noise ratio is a single point in the space this review is asking about, and the record read
here does not identify which DeepStrain release or weights were used, so the result cannot be tied to
a version. Corroboration would require the same grid applied to at least one further estimator from a
different family and to at least one further generator, so that an effect of the estimator can be
separated from an effect of the image-formation model.

What the experiment does establish is mixed, and both halves deserve equal weight. A learning-based
estimator underestimated radial strain across all four cases, most visibly where the wall was thin,
and that is a genuine finding. In the same experiment, on the same images and the same estimator,
circumferential strain was recovered to within an average error of $0.02\pm0.04$, which includes the
infarcted case. The one published test of a cine estimator against a prescribed focal deficit
therefore reports an attribute-specific result rather than a uniform loss, which is the empirical
instance of the claim form set out in \cref{sec:recoverability}: an estimator may hold one component
while losing another. Several explanations for the radial result are available and the single
experiment separates none of them, among them the composition of the training set, which contained
mostly non-infarcted anatomy; wall thinning and partial-volume effects in the scar region; the
generator's image-formation model; and the estimator's own regularisation. Nor is the result a
measurement of how the underestimation scales with scar size, transmurality or location, because
those were not varied.

The other prescribed-motion studies cover other axes but not this one. Adams et al.\ varied spatial
and temporal resolution and showed that peak strain departs further from ground truth as pixel size
grows, with the effect differing between two versions of the same commercial
package~\citep{adams2023}. DENSE-SIM varied the regularisation strength of the analysis tool and
showed that it trades bias against variance in regional radial strain, worst near the
endocardium~\citep{barbaroux2025}. The space-time finite-element framework of Liu et al.\ evaluated
displacement and strain against one synthetic dataset with ground truth and otherwise against
landmarks and masks~\citep{liu2026fe}. STRAUS includes ischaemic cases with prescribed motion fields
and so could in principle support a focal analysis, but the published framework paper reports
discrimination of pathological from healthy subjects by global indices and regional strain curves
rather than a quantified regional recovery error for any estimator~\citep{zhou2018straus}. Taken
together: resolution has been varied against known motion, analysis regularisation has been varied
against known motion, and abnormality has been \emph{represented} in known-motion data. Within the
searches recorded in \cref{sec:methods} and up to the stated cutoff, this review did not identify a
study that varied abnormality size, severity or transmurality against known motion while holding the
estimator and the measurement definition fixed. As with every such statement in this review, that is
a property of the recorded searches rather than an assertion that no such study exists.

\subsection{Regional agreement against material-sensitive references}\label{sec:regionalagreement}

Where a prescribed abnormality is unavailable, the closest in-vivo evidence is segmental agreement
with tagging or DENSE, and here there is a consistent and well-replicated finding: regional
agreement is worse than global agreement, and radial strain is worst of all. The numbers were given
in \cref{sec:material}: regional circumferential ICCs of 0.57--0.77 and regional radial ICCs of
0.05--0.51 against tagging in 61 participants~\citep{militaru2021}; significant biases in every
segmental peak-strain parameter with segmental ICCs of 0.62--0.77 except lateral longitudinal strain
at 0.23 in 42 subjects~\citep{brandt2024}; poorer reliability of basal, apical, torsional and
dyssynchrony measures in 88 participants and 186 paired slices~\citep{wehner2018}; and segmental
repeatability substantially worse than global repeatability in the strain-8
experiment~\citep{bell2025strain8}. Supervision on DENSE improves this without eliminating it:
StrainNet reached a segmental circumferential ICC of 0.75 against DENSE, compared with 0.48 for
commercial feature tracking, in a mixed cohort of 305 participants~\citep{wang2023strainnet}.

Three cautions apply to reading this as evidence about focal recovery. These are agreement
statistics, not error against a known deficit, and the reference carries its own regional bias
(\cref{sec:prescribed}). The cohorts are mixed or healthy rather than selected for focal
abnormality, so a segmental ICC averages over segments that mostly have nothing wrong with them,
which is spectrum bias in its usual measurement sense.
And an ICC is sensitive to the between-segment variance in the sample, so the same estimator can
produce different regional ICCs in populations with different distributions of regional
dysfunction. What the body of evidence does establish is that spatial scale and strain component
change the answer, and that a global agreement statistic cannot be assigned to an individual
segment.

\subsection{Inferring infarction from cine alone}\label{sec:cineonly}

A separate and rapidly growing literature asks whether cine alone can identify infarcted or fibrotic
myocardium. Xu et al.\ learned motion features from non-contrast cine with a recurrent architecture
and optical flow and delineated infarction directly, reporting pixel accuracy of 95.03\%, Dice of
89.87\% and $\kappa=0.91$ against manual delineation of delayed-enhancement images by two
radiologists in 165 cine datasets~\citep{xu2018}. Zhang et al.\ diagnosed chronic infarction from
non-enhanced cine~\citep{zhang2019mi}; Righetti et al.\ classified scar presence by adding
Fourier- and monogenic-signal parametric images to cine~\citep{righetti2024}; a preprint
segments scar from cine using an end-diastole-referenced displacement field, with LGE annotations
transferred into cine coordinates by registration~\citep{yang2025fusion}; and cine radiomics has
been used to predict adverse left ventricular remodelling after ST-segment elevation
infarction~\citep{xin2023}. Deep learning synthetic
strain detects regional wall-motion abnormality against a radiologist consensus reference, with an
AUC of 0.90 in 53 patients and 846 segments~\citep{masutani2023dlss}, and deep learning-based
aligned strain improves detection of LGE-defined fibrosis in Duchenne muscular
dystrophy~\citep{koehler2025alignedstrain}. Virtual native enhancement agrees closely with
LGE-derived scar size and transmurality, but its input includes native T1 mapping and so it operates
under a richer tissue-information regime than cine alone~\citep{zhang2022vne}.

These are substantial results, and they are frequently cited as evidence that cine motion estimates
capture regional dysfunction. They do not establish that. In every case the reference standard is a
tissue map or an expert reading, not a measurement of how far the myocardium moved. A model can
reach a high AUC against LGE by exploiting wall thinning, regional texture, timing or any other
correlate, with or without an accurate displacement field; conversely, an estimator could recover
displacement accurately and still discriminate tissue poorly, because motion and tissue abnormality
are not in one-to-one correspondence. The useful reading is the one the studies themselves support, and it is a
positive one: a routine cine series, with no contrast agent, carries enough information for a learned
model to localise infarction at a level that is clinically interesting, with segment-level agreement
with expert reading at an AUC of 0.90 in one study~\citep{masutani2023dlss}, pixel-level agreement
with delayed-enhancement delineation at a Dice of 0.90 in another~\citep{xu2018}. For a patient in
whom gadolinium is contraindicated, or when a scan must be shortened, that is worth having on its own
terms. What it leaves open is the accuracy of the motion field, which is a different claim and needs
a different reference.

\subsection{Acquisition, reconstruction and processing conditions}\label{sec:conditions}

The last group of evidence concerns the images themselves, and here one distinction does most of the
work: between \emph{precision}, which repeated measurement can establish, and \emph{accuracy}, which
requires a reference. A second distinction, between a global and a regional measurement, is equally
important, because almost all of this literature reports global strain.

Accelerated and deep-learning reconstruction has generally not degraded functional measurements in
the studies retrieved. CineVN, a variational-network reconstruction, was evaluated in 18 healthy
subjects prospectively and 46 patients, comparing volumes and strain with conventional
cine~\citep{vornehm2025cinevn}. REGAIN, trained on 1,616 patients and evaluated in 181 individuals
(126 patients, 55 healthy), gave diagnostic-quality and artefact scores similar to generalised
autocalibrating partially parallel acquisition (GRAPPA) with
Bland--Altman agreement for left ventricular function, volume and strain~\citep{yoon2023regain}. A
deformation-encoding transformer for high-frame-rate cine was developed on 5,840 retrospective
patients across nine centres and tested prospectively in 49 patients and 12 healthy subjects, with
image root-mean-square error and reader preference for temporal smoothness as the
endpoints~\citep{morales2024dent}. A model-based deep-learning reconstruction compared with
compressed-sensing SENSE in ten healthy volunteers at reduction factors of 2, 4, 6 and 8 improved
the \emph{precision} of global circumferential strain at factors 6 and 8, measured as Bland--Altman
agreement with the fully sampled acquisition~\citep{tsuneta2025dlr}. Processing choices can matter
more than reconstruction: Grob et al.\ found that retrospective temporal-resolution interpolation of
compressed-sensing cine significantly underestimated peak longitudinal and circumferential strain
($P<0.01$) while volumes and ejection fraction were similar, and left many strain-rate curves
unusable, in 19 patients and 51 acquisitions~\citep{grob2025}. A generalist foundation model for
accelerated cardiovascular imaging has been described in a preprint that is not peer
reviewed~\citep{wang2026cardiomm}, and CMRxRecon2024 provides a multi-contrast, multi-view raw
$k$-space resource that includes both cine and tagging~\citep{wang2025cmrxrecon2024}.

Three readings follow. First, the common summary that \enquote{deep-learning reconstruction
preserves strain} is supported for global strain under the tested conditions, and is untested for
regional strain: none of the studies above reports segmental accuracy against an independent motion
reference. Second, the endpoints differ in kind: agreement with a conventional acquisition
analysed by the same software is a consistency result, image root-mean-square error is an image
result, and precision relative to a fully sampled acquisition is a precision result. None of them
is accuracy against material motion. Third, a point of vocabulary that is easy to slip on: the
undersampling in these studies is undersampling of \emph{$k$-space}, the spatial-frequency domain of
the image. That is a different frequency domain from the spatial-frequency content of the
\emph{motion field} discussed in \cref{sec:unsupervised}. A reconstruction network alters image
texture, and motion estimation depends on texture, so there is a real mechanism by which
reconstruction could affect a motion estimate, but it is a mechanism that must be demonstrated
through the images, not inferred from a shared word.

\section{Synthesis: where the evidence is, and where the gap is}\label{sec:synthesis}

\subsection{Stating the claim precisely}\label{sec:recoverability}

The preceding sections make it clear that \enquote{can cine recover regional dysfunction?} is not a
single question with a single answer, because the answer depends on what is being recovered, under
what conditions, against what reference and within what tolerance. This review uses \emph{conditional
recoverability} for the resulting claim form. A regional attribute, meaning magnitude, location,
extent or timing, is conditionally recoverable when, under declared acquisition and image-quality
conditions, a declared class of abnormality, a named estimator and version, and a complete
measurement specification in the sense of \cref{tab:estimand}, its error against a stated reference
remains within a prespecified tolerance at a reported coverage. \Cref{fig:attributes} separates the
four attributes from the two validity checks that must accompany them.

Three consequences follow from stating it this way. The claim is attribute-specific: an estimator
may hold the location of a deficit while underestimating its amplitude, which is what the MRXCAT2.0
radial result looks like. It is reference-specific: a tolerance met against DENSE is not
automatically met against prescribed motion, and the reference's own error enters the comparison. And
recoverability is not the same as detectability: detectability asks whether an abnormal case can
be distinguished from a reference condition, recoverability asks whether its magnitude, location,
extent and timing are retained, so a classifier that separates infarcted from normal hearts at high
AUC has not thereby recovered the deficit. This is a reporting framework rather than a validated
biomarker property; this review sets no numerical tolerances and does not estimate a recoverability
surface.

\begin{figure}[!t]
 \centering
 \includegraphics[width=\textwidth,keepaspectratio]{figures/p4_Figure4_error_attributes_mapping_validity.pdf}
 \caption{Regional error attributes and mapping validity are separate checks. \textbf{a--f},
 schematic reference and estimate profiles illustrate differences in magnitude, location, extent,
 spillover beyond a declared evaluation region, temporal magnitude and peak time; the curves are
 non-negative schematics and carry no physiological strain sign or measured value.
 \enquote{Declared support} denotes a prespecified evaluation region rather than strict compact
 support. \textbf{g}, mapping validity is a different question: for $\varphi(X,Y)=(X^2,Y)$ the
 Jacobian determinant is $2X$, so the map reverses orientation for $X<0$, is singular at $X=0$, and
 sends the distinct points $-a$ and $+a$ to the same image $a^2$. This is an analytic example, not a
 measured reference--estimate comparison.}
 \label{fig:attributes}
\end{figure}

The remainder of this section organises the evidence along the three axes on which regional recovery
could plausibly depend, and states for each what is established, what is missing and what design
would settle it.

\subsection{Axis 1: abnormality scale and severity}\label{sec:axis_scale}

\textbf{What is established.} A focal deficit can coexist with a preserved global measurement: the
MRXCAT2.0 infarct simulation had an ejection fraction of 49\% with remote compensation, and its
scar-region circumferential strain was 0.01 against $-0.18$ in remote tissue~\citep{mrxcat2023}. One
estimator evaluated on that simulation underestimated radial strain, with the largest Dice and
strain errors where the wall was thin, while recovering circumferential strain in the same images to
within an average error of $0.02$. Recovery of a focal deficit is therefore attribute-specific in
the one place it has been measured directly. In vivo, segmental agreement with material-sensitive
references is systematically weaker than global agreement, and radial strain is the weakest
component~\citep{militaru2021,brandt2024,wehner2018}; supervision on DENSE narrows but does not
close that gap, reaching a segmental circumferential ICC of 0.75 against 0.48 for commercial feature
tracking in 305 participants~\citep{wang2023strainnet}.

\textbf{What is missing.} The dependence of recovery on the \emph{size}, \emph{severity} and
\emph{transmurality} of the abnormality. Every number above comes from a single abnormality
configuration or from a population average over mostly normal segments. Without varying the
abnormality itself, attenuation that is specific to small or thin lesions cannot be distinguished
from a uniform scale bias.

\textbf{What would settle it.} A prescribed-motion experiment in which abnormality size, severity,
transmural extent and location are varied on a grid under a fixed estimator and a fixed measurement
definition, with a no-abnormality control arm so that an apparent deficit can be separated from
noise amplification, and with error reported per attribute: amplitude in percentage points,
localisation and width in millimetres, timing in milliseconds or cardiac phase, and spillover into
neighbouring regions as a declared quantity. Of the two pathology-bearing prescribed-motion
resources in \cref{tab:datasets}, MRXCAT2.0 states that its scar model could be extended to
arbitrary shapes and property variations, so a grid is within the scope its authors
describe~\citep{mrxcat2023}. STRAUS as released is a fixed set of 18 cases whose infarct locations
and sizes are not exposed as parameters, so generating a grid from it would mean re-running the
electromechanical pipeline that produced it rather than using the published data.

\subsection{Axis 2: motion representation and prior}\label{sec:axis_prior}

\textbf{What is established.} Different priors give different regional answers in the same images:
feature tracking, registration and biomechanically constrained estimators disagree more regionally
than globally~\citep{militaru2021,goetze2026}, and the analysis-side regularisation parameter of a
DENSE pipeline measurably trades bias against variance in regional radial
strain~\citep{barbaroux2025}. The latter is the cleanest existing demonstration that a prior
strength, held under known motion, moves a regional number in a predictable direction.

\textbf{What is missing.} The same experiment on the estimation side for cine. The representation
families now available make this unusually tractable: a band-limited displacement representation has
an explicit, tunable bandwidth~\citep{jia2023fouriernet,jia2026fouriernetplus}, a classical
registration has an explicit $\alpha$ in \cref{eq:registration}, and a biomechanical or
physics-informed estimator has explicit model weights~\citep{qin2023biomech,arratia2023warppinn,liu2026fe}.
Within the recorded searches this review did not identify a study that varied such a control under known
motion and measured the effect on a focal deficit.

\textbf{What would settle it, and what it must not assume.} The experiment is to hold images,
implementation and measurement definition fixed, vary one controllable prior, and evaluate against
prescribed motion. Three distinctions have to be preserved in its design and its interpretation.
First, representation capacity, estimation behaviour and input information are three different
things (\cref{sec:illposed}); a result obtained by changing a bandwidth constrains the first two
jointly unless the input is also controlled. Second, a smooth-looking displacement field is not
evidence of a flattened strain field, because strain is a derivative; the strain field must be
measured, not inferred from the appearance of $\vect{u}$. Third, a Fourier basis is global, so
widening a band does not improve the representation of a chosen region, and any regional claim about
bandwidth has to be demonstrated regionally. A further methodological point: changing a
\emph{training-loss} weight requires matched retraining under controlled splits, whereas changing an
exposed \emph{inference-time} control does not; the two are different experiments and are
distinguished in Supplementary Figure~S1.

\subsection{Axis 3: imaging and processing conditions}\label{sec:axis_conditions}

\textbf{What is established.} Spatial resolution changes peak strain error against known motion, by
up to about 9 percentage points for radial strain between 1.4 and 3.0~mm pixels in one
phantom~\citep{adams2023}. Retrospective temporal interpolation significantly underestimates peak
longitudinal and circumferential strain~\citep{grob2025}. Deep-learning and compressed-sensing
reconstructions have not degraded global strain agreement or precision in the cohorts
tested~\citep{vornehm2025cinevn,yoon2023regain,tsuneta2025dlr}. Through-plane motion degrades
segmental repeatability, particularly basally~\citep{bell2025strain8}.

\textbf{What is missing.} Every one of those results is either global or obtained without a regional
motion reference. The interaction is also untested: whether a reconstruction that preserves global
strain also preserves a focal deficit is a separate question from whether it preserves the global
number, and the answer could differ with lesion size.

\textbf{What would settle it.} Paired acquisitions in the same participants across reconstruction
and sampling conditions, analysed with one locked estimator and one estimand, with a regional
reference where one exists; and, in simulation, controlled degradation applied to images whose
underlying motion is known, so that sampling, reconstruction and estimation error can be separated.
\Cref{tab:designs} sets out the three designs together, and \cref{fig:qc} shows how a prespecified
evaluation and an optional quality-control rule fit around them. \Cref{tab:designs} states evidence
needs rather than a powered protocol: the precision each arm requires follows from the tolerance,
and this review deliberately sets none, so no sample size is given and none should be inferred.

The resources these designs would need already exist in part, and the two that carry pathology or a
human material reference are complementary rather than alternative. A prescribed-motion generator
supplies a deficit of known size in a known place, so an experiment on it answers a question about
\emph{accuracy under prescribed motion}; its limits are that the images are synthesised from the
same deformation the estimator is being asked to recover, which flatters an estimator whose prior
resembles the generator's, and that the released case lists are finite. MRXCAT2.0 covers one scar
geometry but states that its model admits others; STRAUS covers 12 ischaemic cases across four
coronary configurations but as a fixed set.

A paired human resource such as the Stanford collection
supplies cine and xyz-encoded DENSE at identical slice locations in the same participants, so an
experiment on it answers a question about \emph{agreement with DENSE}, with the reference's own
processing uncertainty and the gating difference inside the comparison; what it cannot answer is
whether a focal deficit survives, because it contains none. Running both is therefore not redundant:
the simulation bounds the error against a known deficit, the human data bound the transfer of that
estimator to real images and a real reference, and neither substitutes for the other. The prior
sweep of \cref{sec:axis_prior} needs only an estimator with an exposed control and the same
prescribed-motion data. What none of these resources yet provides is a patient cohort with focal
disease \emph{and} a material-sensitive motion reference, which is why \cref{tab:open} lists it as a
field-level gap rather than a study design.

\begin{longtable}{@{}L{0.18\textwidth}L{0.26\textwidth}L{0.28\textwidth}L{0.22\textwidth}@{}}
\caption{Study designs that would close the three gaps. The rows separate evidence needs; they are
not a validated protocol or an accuracy hierarchy.}\label{tab:designs}\\
\toprule
Axis and question & Reference and design & Informative outputs & Interpretation boundary\\
\midrule
\endfirsthead
\toprule
Axis and question & Reference and design & Informative outputs & Interpretation boundary\\
\midrule
\endhead
Abnormality scale and severity: how does recovery change with the size, severity and transmurality of the deficit? & Prescribed motion from an independent generator; abnormality parameters varied on a grid; a no-abnormality control arm; one fixed estimator and estimand & Amplitude error in percentage points; localisation and width in millimetres; timing error in milliseconds or phase; spillover into neighbouring regions; topology checks & Recovery within the simulated regime; not an in-vivo guarantee, and not transferable to another generator without testing\\
Motion representation and prior: which part of a regional error is the prior? & Images, implementation and estimand held fixed; one controllable prior varied (regularisation weight, representation bandwidth, physics weight); matched retraining where the control is a training weight & Attribute-specific error as a function of the control; separation of representation capacity, estimation behaviour and input information & Constrains the varied control only; a monotone trend in one estimator is not a general law of regularisation\\
Imaging and processing conditions: do sampling and reconstruction change the regional answer? & Paired acquisitions in the same participants across reconstruction and sampling conditions with a locked estimator; plus controlled degradation of images with known motion & Regional error and precision separately; interaction between condition and abnormality size; failure and exclusion rates & Precision results do not establish accuracy; $k$-space sampling and motion-field frequency are different quantities\\
External transfer: do regional findings and failure indicators carry over? & Held-out anatomies and sites; paired cine with DENSE or tagging; repeat cine; multi-scanner data; LGE for tissue stratification & Agreement against a registered material-sensitive reference with its uncertainty; repeatability; subgroup performance; calibration of any failure rule & Reference-specific; limited samples and imperfect references narrow the claim\\
Decision contribution: does a specified measurement add useful information? & Fixed estimator, estimand and failure handling; cohort with ejection fraction, infarct variables and outcomes; held-out evaluation & Incremental calibration and discrimination; decision analysis; threshold transfer; profile of accepted and rejected cases & Association does not establish management benefit without an appropriate intervention study\\
\bottomrule
\end{longtable}

\begin{figure}[!t]
 \centering
 \includegraphics[width=\textwidth,keepaspectratio]{figures/p5_Figure5_prespecified_evaluation_QC.pdf}
 \caption{A prespecified evaluation, with an optional quality-control or abstention rule.
 \textbf{a}, acquisition and image quality, abnormality characteristics, and estimator and version
 are declared under a fixed measurement definition. \textbf{b}, the comparison names its evaluation
 reference and that reference's uncertainty, an attribute-specific loss $L_g$ and a prespecified
 tolerance $\delta_g$. \textbf{c}, held-out cases are evaluated for magnitude, location, extent and
 timing, with separate implementation and field-validity checks. \textbf{d}, any reporting or
 abstention rule is fixed before final testing and operates on observable quality-control signals,
 not on the unknown true error; the joint audit then reports coverage together with accepted-case
 errors, failures and subgroups. This is a proposed evaluation design, not a validated
 quality-control threshold or a clinical decision rule.}
 \label{fig:qc}
\end{figure}

\subsection{Keeping the evidence layers apart}\label{sec:layers}

The three axes concern measurement. A separate hierarchy governs what a measurement licenses
clinically, and conflating its levels is the most common way that regional claims become
overstated. Technical accuracy against a reference, precision under repeated measurement,
association with tissue or outcome, incremental value over existing covariates, and benefit from
acting on a threshold are five distinct claims requiring five distinct study
designs~\citep{raunig2015,fda2022qib,puntmann2018}. Ejection fraction illustrates the complete chain
rather than the absence of error: it has a chamber-volume definition, standardised reporting and
guideline categories tied to intervention evidence~\citep{schulzmenger2020,hundley2022,heidenreich2022},
even though core-laboratory comparisons found that fewer than 55\% of paired measurements agreed
within five percentage points across the tested modalities~\citep{pellikka2018}. Strain does not yet
have that chain, and the gap is more than technical: a threshold must name a quantity, a population
and a decision, and be connected to evidence that the resulting action helps. Echocardiographic
consensus recommendations provide a precedent for specifying acquisition, analysis, reporting and
intended use~\citep{thomas2025}, and CMR reference values remain conditional on technique,
population and software~\citep{kawelboehm2025}. A version-locked workflow can be useful in a
specialist setting well before any of this is settled; the appropriate claim is simply local.

\section{Other open problems and conclusion}\label{sec:open}

\subsection{Open problems beyond the three axes}\label{sec:otheropen}

The three axes of \cref{sec:synthesis} follow from the central question. Several further problems
recur across the literature reviewed and are set out in \cref{tab:open} because they condition any
answer to it.

\begin{longtable}{@{}L{0.20\textwidth}L{0.33\textwidth}L{0.42\textwidth}@{}}
\caption{Further open problems in regional cine motion and strain estimation.}\label{tab:open}\\
\toprule
Problem & Where the evidence currently stands & What would move it forward\\
\midrule
\endfirsthead
\toprule
Problem & Where the evidence currently stands & What would move it forward\\
\midrule
\endhead
Through-plane motion and 3D from 2D & Multi-view and mesh-based estimators now produce 3D fields from 2D acquisitions~\citep{meng2022mulvimotion,meng2024deepmesh,liu2026fe}, but their published evaluation is mask or landmark agreement, and segmental repeatability remains worst basally~\citep{bell2025strain8} & Simulations with prescribed through-plane motion; comparison against 3D tagging or DENSE; reporting of behaviour when one view is missing or misaligned\\
Generalisation across acquisition and population & Learning-based estimators are predominantly developed on UK Biobank, ACDC and M\&Ms; online adaptation has been proposed for distribution shift~\citep{yu2020foal} & Error reported stratified by disease class, field strength, vendor and reconstruction, rather than pooled\\
Uncertainty and failure detection & Bayesian and probabilistic estimators now output uncertainty maps~\citep{chen2022transmorph,movahed2026cardiomorphnet} & Calibration of the uncertainty against measured regional error on data with a reference; coverage and accepted-case error reported together, following the selective-prediction formulation~\citep{geifman2019selectivenet}\\
Serial precision and minimum detectable change & Global coefficients of variation are available~\citep{bell2025strain8}, and DeepStrain reported scan--rescan in ten healthy volunteers~\citep{morales2021} & Segment-level scan--rescan studies in patient populations, reported as a repeatability coefficient with its confidence limits and distinguished from a clinically important change\\
A public cine benchmark with a motion reference & CMAC supplies 12 landmarks in 15 healthy volunteers~\citep{tobongomez2013}; the Stanford collection supplies cine and xyz-encoded DENSE at identical slice locations, plus mid-ventricular tagging and native T1, in 55 participants of whom 51 are healthy and 4 have aortic stenosis~\citep{stanforddata}; synthetic resources supply prescribed motion~\citep{mrxcat2023,zhou2018straus,barbaroux2025} & An open benchmark that pairs cine with a material-sensitive reference in \emph{patients with focal disease}, or a realistic simulation resource with per-case reference motion and a declared abnormality grid\\
Cross-modality precedent & In echocardiography, a photorealistic simulation strategy incorporating speckle decorrelation from real recordings produced an open dataset of 1,478 videos with reference motion, used to train a motion estimator reported to reach a global longitudinal strain variability of 1.42\% against 1.78\% for the clinical reference~\citep{judge2026} (preprint; speckle-tracking echocardiography, not CMR) & The analogous resource does not exist for cine CMR. Whether a cine image-formation model can reach comparable realism is itself an open question\\
\bottomrule
\end{longtable}

\subsection{Limitations of this review}\label{sec:limitations}

This is a critical narrative review, not a systematic review or a pooled accuracy assessment.
Retrieval was limited by indexing, query vocabulary and result caps, and screening was performed by
one person without independent duplicate screening; no claim of exhaustive retrieval is made, and
statements that something was \enquote{not identified} are bounded by the recorded searches and the
1 September 2026 cutoff rather than being assertions of absence. No screening flow is reported: the
query strings, database caps and hit counts are supplied in the supplementary search record, but no
count of records screened, excluded or included is given, so a reader cannot distinguish a wide
search with narrow inclusion from a narrow search. The review was not registered and no protocol was
deposited.

A further limit applies to the central observation that deliberate evaluation on focal
abnormality is rare. The observation is made on the published record, and a published record can be
thin for reasons other than the experiments not having been performed, including page limits in
conference and workshop venues and the ordinary incentive to report the evaluation that supports a
method's stated claim. This bounds the inference, and it is not a criticism of any cited author.

Heterogeneous populations,
acquisitions, software versions, references and measurement definitions preclude a shared error
distribution, which is why no result in \cref{tab:methods_validation} has been pooled or ranked.
Several sources were read at abstract level only; those are marked in the accompanying claim audit,
and no full-text-level statement is made from them. Bibliographic verification confirms that a
record exists and says what it is; it does not validate the study. Conditional recoverability is the
interpretive frame adopted here, not a construct that the cited studies set out to test, and this
review supplies no numerical tolerances, no recoverability surface and no validated quality-control
rule. The review is also restricted to ischaemic regional dysfunction as its running example; a 2026
scoping review synthesising strain in non-ischaemic cardiomyopathy covers both echocardiography and
CMR, and its aggregate conclusions cannot all be assigned to routine cine~\citep{mohammed2026}.
Finally, this review is restricted to the left ventricle and principally to short-axis analysis; right ventricular, atrial and layer-specific measurement raise the same questions under
harder conditions.

\subsection{Conclusion}\label{sec:conclusion}

Routine cine CMR supports motion and strain estimation without dedicated motion encoding, but dense
regional material correspondence is estimated rather than observed, and that single fact organises
the evidence. Returning to the question posed in \cref{sec:introduction}: current evidence supports
regional recovery \emph{conditionally} and \emph{unevenly}. For global measurements, the evidence is
substantial, because repeatability, cross-method agreement and prognostic association are all
documented.
For regional measurements, agreement with material-sensitive references is consistently weaker than
global agreement and weakest for radial strain, which is a well-replicated finding across
independent cohorts. For focal abnormality specifically, the direct evidence is thin and mixed
rather than thin and negative: within the literature retrieved, systematic evaluation against a
prescribed focal deficit appears in a single independent synthetic study with one scar geometry and
one estimator, and in that study circumferential strain was recovered to within an average error of
$0.02$ while radial strain was underestimated by about $0.24$. The large and growing literature that
identifies infarction from cine validates against tissue references, which answer a different
question.

That is a statement about the evidence, not about the methods. Nothing reviewed here establishes
that cine-based estimators fail to recover focal deficits; the point is that the experiments that
would show either way, namely varying abnormality size and severity under known motion, varying a
controllable prior while holding images and estimand fixed, and varying acquisition and
reconstruction conditions with a regional reference, were not identified within the searches
recorded in \cref{sec:methods}. The methods
literature has moved quickly towards three-dimensional, temporally continuous and physically
constrained estimators, and the validation literature has not moved with it. Until it does, a
regional claim should name its implementation, its measurement definition, its population, its
reference, its tolerance and its coverage, and should be read as local to them.
